% !TeX root = ../../Book4.tex
% 纲要标题：## 附录 A　单子、下降与可表性
% 纲要位置：工作记录/计划纲要.md，第 3816 行。
% 纲要细目（写作范围）：
% > 完整构造单子代数、自由伴随及比较函子；证明分裂余等化版本的 Barr–Beck 判据、忠实平坦模下降及带解集条件的一般伴随函子定理。


\chapter{单子、下降与可表性}
\label{b4:app:A}

\begin{chapterabstract}
本附录研究伴随产生的单子与其代数、比较函子和单子性判据，并以模下降检验恢复对象的条件。解集条件与极限保持给出伴随存在的判据；Kleisli 范畴\footnote{克莱斯利（Heinrich Kleisli，1930-2011），瑞士数学家，主要研究代数拓扑与范畴论。}与 Lawvere 代数理论\footnote{劳维尔（Francis William Lawvere，1937-2023），美国数学家，主要以范畴语言研究代数理论和数学基础。}则分别把形式替换和有限变量运算组织成范畴。
\end{chapterabstract}

\begin{learninggoals}
\begin{itemize}
\item 按单位、乘法及求值等式识别单子代数，构造自由代数与其伴随。
\item 区分比较函子的存在、保守性和单子性，逐项核验分裂余等化子条件。
\item 用余单子写出模下降数据，并从相容元素构造恢复模。
\item 把可表性与伴随存在性化为始对象问题，辨认极限假设和大小假设的不同作用。
\item 计算 Kleisli 复合，构造其伴随，并把它的态射识别为自由代数之间的同态。
\item 用有限积记录运算项的代入，从群和模的代数理论恢复模型及其同态。
\end{itemize}
预备知识为局部小范畴、自然变换、伴随、等化子与余等化子，以及模的张量积和正合性。关于范畴大小，所有“小”均相对于固定的集合范围；不对真类大小的对象族任意取积。
\end{learninggoals}

设 \(X\) 是一个集合。由 \(X\) 的字组成自由幺半群 \(T(X)\)；若把这些字再排成一个字，就得到 \(T(T(X))\)。把内层的字依次接起来，可以从 \(T(T(X))\) 回到 \(T(X)\)。连续做三次自由构造时，先压平内两层还是先压平外两层，结果都应是同一个长字。这个可以直接用拼接核对的事实，提示了单子乘法的结合律。对任意代数，类似的「先代入再求值」为什么也应相容？从这个问题出发，我们将比较自由构造、单子代数与原来的对象范畴。

\ProseParagraphBreak
正文中的伴随经常来自已经知道的自由对象。这里还要追问两件事：伴随产生的单子是否保存了另一范畴的全部信息，以及在尚未构造自由对象时，哪些条件能保证它存在。前一个问题由单子性判据处理，后一个问题由伴随函子定理处理。下降则提供一个具体用途：从扩张后的模及其相容数据恢复原模。伴随数据的基本定义可对照 Weibel 的叙述\cite[A.6.1; Theorem A.6.2, pp. 429--431]{Weibel1994HomologicalAlgebra}。

\ProseParagraphBreak
\chapref{b4:ch:03}的伴随与单子为本附录提供基础。单子代数与自由代数使“归并”成为可以求值的结构；比较函子与单子性判据进一步检验原来的范畴能否从这一结构恢复，模下降给出具体的恢复计算。可表性又把伴随存在的问题化为始对象问题，需要同时检查极限条件与解集条件。最后，把形式表达式之间的替换视为态射，就得到 Kleisli 范畴；仅记录有限变量及其运算项，还能用有限积描述群、模等代数结构。

% !TeX root = ../../Book4.tex
% 纲要标题：### A.1 单子与余单子
% 纲要位置：工作记录/计划纲要.md，第 3820 行。
% 纲要细目（写作范围）：
% #### A.1.1 单子的单位与乘法
% #### A.1.2 余单子与伴随所给的构造
% #### A.1.3 幂等单子与反射子范畴


\section{单子与余单子}
\label{b4:sec:A01}

一个自由对象的元素往往是形式表达式。把表达式再放进同一种形式构造，会出现两层括号；单子的乘法负责把它们合并，单位则把原对象放进形式构造。伴随所给的结构满足明确的单位与结合等式；这些等式也有助于分辨单子与反射之间的关系。

\subsection{单子的单位与乘法}
\label{b4:subsec:A01:01}

沿用\cref{b4:def:monad}，范畴 \(\mathcal C\) 上的单子为自函子 \(T\) 及自然变换
\(\eta:1_{\mathcal C}\Rightarrow T\)、\(\mu:T^2\Rightarrow T\)，满足
\begin{equation}\label{b4:eq:A-monad-laws}
\mu_XT(\mu_X)=\mu_X\mu_{TX},\qquad
\mu_XT(\eta_X)=1_{TX}=\mu_X\eta_{TX}.
\end{equation}
这里及本附录中，态射乘积按从右到左的顺序复合。例如第一式的两边都从 \(T^3X\) 到 \(TX\)，第二式的两边都从 \(TX\) 到自身。

\ProseParagraphBreak

在集合范畴上，令 \(TX\) 为由 \(X\) 中元素组成的全部有限有序表，包括空表。单位把 \(x\) 送到单元素表 \([x]\)，乘法把一个“表的表”按顺序连接。把三层表连接为一层时，先连接内层或先连接外层结果相同；加入单元素表不改变内容。这逐项验证了\eqref{b4:eq:A-monad-laws}。这里的表保留次序，不能把它擅自换成无序的有限集合。

\begin{example}\label{b4:exm:A-extension-scalars-monad}
设 \(\varphi:R\rightarrow S\) 为含幺环同态，在左 \(R\)-模范畴上令
\[
T(M)=S\otimes_RM,\qquad
\eta_M(m)=1\otimes m,\qquad
\mu_M(s\otimes t\otimes m)=st\otimes m.
\]
构造张量积时使用 \(S\) 的右 \(R\)-作用，而所得对象的左 \(R\)-作用来自 \(S\) 的左乘法：
\[
s\cdot r=s\varphi(r),\qquad
s\varphi(r)\otimes m=s\otimes rm,\qquad
r\cdot(s\otimes m)=\varphi(r)s\otimes m.
\]
在非交换情形，\(\varphi(r)s\otimes m\) 与 \(s\varphi(r)\otimes m\) 一般不同；构造依靠的是 \(S\) 的左 \(S\)-作用与右 \(R\)-作用相容，并未交换标量的次序。两层张量积的平衡关系分别把
\(s\varphi(r)\otimes t\otimes m\) 与 \(s\otimes\varphi(r)t\otimes m\) 识别，以及把
\(s\otimes t\varphi(r)\otimes m\) 与 \(s\otimes t\otimes rm\) 识别。它们在 \(\mu_M\) 下分别有相同的像，所以 \(\mu_M\) 良定义。其左 \(R\)-线性由
\[
\mu_M\bigl(\varphi(r)s\otimes t\otimes m\bigr)
=\varphi(r)st\otimes m
=r\cdot\mu_M(s\otimes t\otimes m)
\]
给出。环乘法的结合律给出单子结合律，\(S\) 的单位元给出两条单位律。因此这是一个单子，不要求 \(R,S\) 交换或 \(S\) 在 \(R\) 上平坦。
\end{example}

\subsection{余单子与伴随所给的构造}
\label{b4:subsec:A01:02}

余单子反转上面的箭头。\cref{b4:def:comonad}中的数据为
\(Q:\mathcal D\to\mathcal D\)、\(\varepsilon:Q\Rightarrow1_{\mathcal D}\)、
\(\delta:Q\Rightarrow Q^2\)，条件为
\[
Q(\delta_Y)\delta_Y=\delta_{QY}\delta_Y,\qquad
Q(\varepsilon_Y)\delta_Y=1_{QY}=\varepsilon_{QY}\delta_Y.
\]
乘法把两层归并，余乘法则把一层放进两层，而且两种继续展开的方法必须一致。

\ProseParagraphBreak

若 \(F:\mathcal C\to\mathcal D\) 与 \(U:\mathcal D\to\mathcal C\) 满足 \(F\dashv U\)，其单位、余单位为 \(\eta,\varepsilon\)，则
\[
T=UF,\quad \mu_X=U(\varepsilon_{FX}),\qquad
Q=FU,\quad \delta_Y=F(\eta_{UY}).
\]
\cref{b4:prop:adjunction-monad}已通过自然性及三角恒等式证明第一组数据是单子；在对偶范畴中应用同一个证明，就得到第二组余单子。本附录固定以 \(U\) 表示右伴随，避免把原对象与构造出的代数遗忘函子混淆。

\begin{definition}\label{b4:def:A-comonad-coalgebra}
设 \(\mathcal D\) 为范畴，\((Q,\varepsilon,\delta)\) 为其上的余单子，\(Y\in\mathcal D\)，\(c:Y\to QY\) 为态射。若满足
\[
\varepsilon_Yc=1_Y,\qquad \delta_Yc=Q(c)c,
\]
则称 \((Y,c)\) 为 \(Q\) 上的\term[yudanzi-shang-yudaishu]{余代数}[coalgebra for a comonad]。对两个余代数 \((Y,c)\)、\((Y',c')\)，若 \(\mathcal D\) 中的态射 \(f:Y\to Y'\) 满足 \(Q(f)c=c'f\)，则称 \(f\) 为它们之间的余代数同态。这些对象与同态组成的范畴记为 \(\mathcal D_Q\)。
\end{definition}
\symidx{\(\mathcal D_Q\)：余单子上的余代数范畴}

例如，固定集合 \(A\)，对余单子 \(QX=A\times X\)，余单位为第二投影，余乘法为 \((a,x)\mapsto\bigl(a,(a,x)\bigr)\)。其余代数恰好是一个带映射 \(r:X\to A\) 的集合：条件迫使 \(c(x)=\bigl(r(x),x\bigr)\)，而第二条公理随后自动成立。同态就是保持到 \(A\) 的映射的函数，所以 \(\mathbf{Set}_Q\) 是切片范畴 \(\mathbf{Set}/A\)。

\subsection{幂等单子与反射子范畴}
\label{b4:subsec:A01:03}

\begin{definition}\label{b4:def:A-idempotent-monad}
设 \((T,\eta,\mu)\) 为范畴 \(\mathcal C\) 上的单子。如果每个乘法态射 \(\mu_X:T^2X\to TX\) 都为同构，则称这个单子\term[mideng-danzi]{幂等}[idempotent monad]。
\end{definition}

\begin{proposition}\label{b4:prop:A-idempotent-reflection}
设 \((T,\eta,\mu)\) 为局部小范畴 \(\mathcal C\) 上的幂等单子。令 \(\mathcal C_{\mathrm{fix}}\) 为满足 \(\eta_Y:Y\to TY\) 可逆的对象 \(Y\in\mathcal C\) 所成的全子范畴。则对每个 \(X\in\mathcal C\)，\(TX\) 都在其中，且对 \(Y\in\mathcal C_{\mathrm{fix}}\)，预复合 \(\eta_X\) 给出自然双射
\[
\operatorname{Hom}_{\mathcal C}(TX,Y)
\longrightarrow\operatorname{Hom}_{\mathcal C}(X,Y).
\]
因此 \(T:\mathcal C\to\mathcal C_{\mathrm{fix}}\) 是包含函子的左伴随。
\end{proposition}
\begin{proof}
两条单位律及 \(\mu_X\) 可逆给出
\(T(\eta_X)=\eta_{TX}=\mu_X^{-1}\)，所以 \(TX\) 属于所述子范畴。
给定 \(f:X\to Y\)，候选延拓为 \(\overline f=\eta_Y^{-1}T(f)\)。由 \(\eta\) 的自然性，\(\overline f\eta_X=f\)。

若 \(h:TX\to Y\) 满足 \(h\eta_X=f\)，则
\[
\eta_Yh=T(h)\eta_{TX}
=T(h)T(\eta_X)=T(f),
\]
从而 \(h=\overline f\)。逆公式由函子及自然变换的复合组成，故对 \(X,Y\) 自然。
\end{proof}

这里可逆的是 \(\mu_X\) 和 \(\eta_{TX}\)，不能据此断定第一次施加 \(T\) 时的 \(\eta_X\) 也可逆。例如在交换群范畴上取有理化单子
\[
T(A)=\mathbb Q\otimes_{\mathbb Z}A.
\]
乘法 \(\mathbb Q\otimes_{\mathbb Z}\mathbb Q\longrightarrow\mathbb Q\) 是同构，其逆为 \(q\mapsto q\otimes1\)，所以 \(T^2\cong T\)。但在 \(A=\mathbb Z\) 处，单位是通常的包含映射 \(\mathbb Z\hookrightarrow\mathbb Q\)，并非同构。幂等性说明再次有理化不再改变对象；满足 \(\eta_Y\) 可逆的固定对象只组成原范畴的一个全子范畴。

\ProseParagraphBreak

全子范畴的包含函子若有左伴随，该子范畴称为\term[fanshe-zifanchou]{反射子范畴}[reflective subcategory]。例如交换化 \(G\mapsto G/[G,G]\) 再作用一次不再改变对象；\cref{b4:exm:abelianization-adjunction}的伴随因而给出幂等单子，其固定对象正是交换群。一般单子却不幂等：自由幺半群中的一个表，通常有许多不同的分段方式，它们在乘法下都连接成同一个表。

% !TeX root = ../../Book4.tex
% 纲要标题：### A.2 单子上的代数与自由对象
% 纲要位置：工作记录/计划纲要.md，第 3828 行。
% 纲要细目（写作范围）：
% #### A.2.1 单子上的代数与同态
% #### A.2.2 自由代数与遗忘伴随
% #### A.2.3 比较函子与典型代数范畴


\section{单子上的代数与自由对象}
\label{b4:sec:A02}

单子给出形式运算，但形式表达式还须在一个对象中求值。求值要把单个元素还原为自身，并要求“先在内层求值再在外层求值”与“一次归并后求值”相同。单子上的代数就是满足这两个条件的对象。

\subsection{单子上的代数与同态}
\label{b4:subsec:A02:01}

\begin{definition}\label{b4:def:A-monad-algebra}
设 \((T,\eta,\mu)\) 为局部小范畴 \(\mathcal C\) 上的单子，\(X\in\mathcal C\)，\(a:TX\to X\) 为态射。若满足
\begin{equation}\label{b4:eq:A-algebra-laws}
a\eta_X=1_X,\qquad aT(a)=a\mu_X,
\end{equation}
则称 \((X,a)\) 为 \(T\) 上的\term[danzi-shang-daishu]{代数}[algebra for a monad]。对两个 \(T\)-代数 \((X,a)\)、\((Y,b)\)，若 \(\mathcal C\) 中的态射 \(f:X\to Y\) 满足
\[
fa=bT(f)
\]
则称 \(f\) 为它们之间的代数同态。所有这些代数与同态组成的范畴记为 \(\mathcal C^T\)，称为 \term[eilenberg-moore-fanchou]{Eilenberg–Moore 范畴}[Eilenberg--Moore category]。\footnote{摩尔（John Coleman Moore，1923-2016），美国数学家，研究代数拓扑和谱序列；艾伦伯格已见本卷前注。}
\end{definition}
\symidx{\(\mathcal C^T\)：单子 \(T\) 上的代数范畴}

恒等态射满足同态等式；若 \(f a=bT(f)\)、\(g b=cT(g)\)，则
\((gf)a=cT(gf)\)。所以它们确实组成范畴。遗忘函子
\(U^T:\mathcal C^T\to\mathcal C\) 把 \((X,a)\) 送到 \(X\)，把同态送到底层态射；它是忠实函子，但一般不是全函子。

\begin{proposition}\label{b4:prop:A-idempotent-algebras}
设 \((T,\eta,\mu)\) 为局部小范畴 \(\mathcal C\) 上的幂等单子，\(X\in\mathcal C\)。\(X\) 存在 \(T\)-代数结构，当且仅当单位分量 \(\eta_X:X\to TX\) 可逆；此时结构唯一，为 \(a=\eta_X^{-1}\)。满足此条件的对象之间的每个 \(\mathcal C\)-态射都是代数同态。
\end{proposition}
\begin{proof}
若 \(a\eta_X=1_X\)，由 \(\eta\) 的自然性及幂等性得到
\[
\eta_Xa=T(a)\eta_{TX}=T(a)T(\eta_X)=1_{TX}.
\]
所以 \(a\) 必为 \(\eta_X\) 的逆。反之，若 \(\eta_X\) 可逆，令 \(a=\eta_X^{-1}\)。单位律显然成立，且
\[
T(a)=T(\eta_X)^{-1}=\mu_X,
\]
从而 \(aT(a)=a\mu_X\)。若 \(f:X\to Y\)，自然性 \(T(f)\eta_X=\eta_Yf\) 在两边乘以单位的逆，就得到代数同态等式。
\end{proof}

\subsection{自由代数与遗忘伴随}
\label{b4:subsec:A02:02}

\begin{theorem}[自由代数与遗忘伴随]\label{b4:thm:A-em-adjunction}
设 \((T,\eta,\mu)\) 为局部小范畴 \(\mathcal C\) 上的单子，\(\mathcal C^T\) 为其代数范畴，\(U^T:\mathcal C^T\to\mathcal C\) 为取底层对象与态射的遗忘函子。对 \(X\in\mathcal C\) 及 \(\mathcal C\) 中的态射 \(f:X\to X'\)，公式
\[
F^T(X)=(TX,\mu_X),\qquad F^T(f)=T(f)
\]
定义函子 \(F^T:\mathcal C\to\mathcal C^T\)，且 \(F^T\dashv U^T\)。对任意 \(T\)-代数 \((Y,b)\)，伴随双射为
\begin{equation}\label{b4:eq:A-em-adjunction}
\operatorname{Hom}_{\mathcal C^T}\bigl(F^TX,(Y,b)\bigr)
\cong\operatorname{Hom}_{\mathcal C}(X,Y),
\qquad h\longmapsto h\eta_X,
\end{equation}
其逆把 \(f\) 送到 \(bT(f)\)。这个伴随产生的单子正是原来的 \(T\)。
\end{theorem}
\begin{proof}
\eqref{b4:eq:A-monad-laws}说明 \((TX,\mu_X)\) 满足代数公理。乘法的自然性
\(T(f)\mu_X=\mu_YT^2(f)\) 说明 \(T(f)\) 是代数同态，函子性来自 \(T\)。

对 \(f:X\to Y\)，令 \(\overline f=bT(f)\)。利用代数公理及乘法自然性，
\[
\overline f\mu_X
=b\mu_YT^2(f)
=bT(b)T^2(f)=bT(\overline f),
\]
所以它是代数同态。又
\(\overline f\eta_X=b\eta_Yf=f\)。
反过来，若 \(h:(TX,\mu_X)\to(Y,b)\) 为代数同态，则
\[
bT(h\eta_X)=bT(h)T(\eta_X)
=h\mu_XT(\eta_X)=h.
\]
故两公式互逆。与源、目标的态射复合时，这两个公式仍相容，因而构成自然双射。

单位为 \(\eta_X\)，余单位在 \((Y,b)\) 处就是 \(b:(TY,\mu_Y)\to(Y,b)\)。所以遗忘后复合自由函子仍为 \(T\)，伴随产生的乘法在 \(X\) 处为自由代数的结构态射 \(\mu_X\)，与原数据相同。
\end{proof}

例如，有限表单子的代数给出幺半群。对求值 \(a:TX\to X\)，定义
\[
e=a\bigl([]\bigr),\qquad x\cdot y=a\bigl([x,y]\bigr).
\]
对由三张单元素表组成的分组，以及把其中两张先连接的分组，代数结合律说明
\((x\cdot y)\cdot z=x\cdot(y\cdot z)=a\bigl([x,y,z]\bigr)\)；对含空表的分组得到左右单位律。归纳可知 \(a\bigl([x_1,\ldots,x_n]\bigr)=x_1\cdots x_n\)。反过来，幺半群的有序乘积满足两条代数公理。代数同态恰好保持空表和二元素表的值，也就是幺半群同态。

\ProseParagraphBreak

对任意 \(T\)-代数 \((Y,b)\)，单位律 \(b\eta_Y=1_Y\) 说明余单位 \(b\) 在底层范畴 \(\mathcal C\) 中是分裂满态射，但这并不自动给出代数范畴中的分裂。有限表单子把幺半群 \(Y\) 的表送到乘积，底层截面为 \(y\mapsto[y]\)；然而
\[
\eta_Y(y_1y_2)=[y_1y_2],\qquad
\eta_Y(y_1)\eta_Y(y_2)=[y_1,y_2],
\]
两张表不同，所以这个截面不是幺半群同态。某些求值映射甚至不存在任何幺半群同态截面：取二元群 \(Y=\{e,g\}\)，其中 \(g^2=e\)。若有幺半群同态 \(\sigma:Y\rightarrow TY\)，则 \(\sigma(g)\sigma(g)=\sigma(e)=[]\)，从表长相加可知 \(\sigma(g)=[]\)。于是 \(b\sigma(g)=e\ne g\)，故 \(\sigma\) 不可能是 \(b\) 的截面。要从自由代数恢复一般代数，还需把求值所满足的关系纳入构造；下一节将用余等化子完成这一步。

\subsection{比较函子与典型代数范畴}
\label{b4:subsec:A02:03}

\begin{proposition}\label{b4:prop:A-comparison-functor}
设 \(F:\mathcal C\to\mathcal D\)、\(U:\mathcal D\to\mathcal C\) 为局部小范畴之间的伴随 \(F\dashv U\)，单位、余单位为 \(\eta,\varepsilon\)，所生单子为 \(T=UF\)。则
\[
K(Y)=\bigl(UY,U(\varepsilon_Y)\bigr),\qquad K(f)=U(f)
\]
定义函子 \(K:\mathcal D\to\mathcal C^T\)，满足 \(U^TK=U\)、\(KF=F^T\)。它称为该伴随的\term[bijiao-hanzi]{比较函子}[comparison functor]。
\end{proposition}
\begin{proof}
记 \(a_Y=U(\varepsilon_Y)\)。三角恒等式给出 \(a_Y\eta_{UY}=1_{UY}\)。把余单位自然性用于 \(\varepsilon_Y:FUY\to Y\)，得到
\[
\varepsilon_YFU(\varepsilon_Y)
=\varepsilon_Y\varepsilon_{FUY}.
\]
应用 \(U\)，便有 \(a_YT(a_Y)=a_Y\mu_{UY}\)，所以 \(K(Y)\) 是代数。对 \(f:Y\to Z\)，余单位自然性又给出
\(U(f)a_Y=a_ZT\bigl(U(f)\bigr)\)，因此 \(K(f)\) 是代数同态。其余函子等式来自定义。
\end{proof}

等式 \(KF=F^T\) 和 \(U^TK=U\) 表明，原伴随的自由对象与单子重新构造的自由代数总能对应，且比较函子保持底层对象与态射。问题是原范畴 \(\mathcal D\) 中的全部对象和态射能否由这些代数数据恢复。为此需要检查 \(K\) 是否为范畴等价：每个 \(T\)-代数都应同构于某个 \(K(Y)\)，而 \(Y,Z\) 之间的态射也应与 \(K(Y),K(Z)\) 之间的代数同态一一对应。仅有自由伴随 \(F^T\dashv U^T\) 的存在还不能保证这两件事。

\begin{definition}\label{b4:def:A-monadic-functor}
设 \(\mathcal C,\mathcal D\) 为局部小范畴，函子 \(U:\mathcal D\to\mathcal C\) 有左伴随 \(F:\mathcal C\to\mathcal D\)，伴随的余单位为 \(\varepsilon:FU\Rightarrow1_{\mathcal D}\)。以 \(\mathcal C^{UF}\) 表示所生单子的代数范畴，定义比较函子 \(K(Y)=(UY,U(\varepsilon_Y))\)、\(K(f)=U(f)\)。若 \(K:\mathcal D\to\mathcal C^{UF}\) 是范畴等价，则称 \(U\) 为\term[danzixing-hanzi]{单子性函子}[monadic functor]，也说这个伴随具有单子性。
\end{definition}

对自由幺半群与底集合的伴随，比较函子把幺半群送到“有限表取乘积”的代数。上一小节已经证明每个代数都唯一确定这种乘法，而且代数同态就是幺半群同态，所以遗忘函子具有单子性。\(\mathcal C^T\) 在这里保留了全部代数结构。

\ProseParagraphBreak

并非每个伴随都能如此恢复另一范畴。离散拓扑函子
\(F:\mathbf{Set}\to\mathbf{Top}\) 左伴随于底集合函子 \(U\)，且 \(UF=1_{\mathbf{Set}}\)，所生单子只是恒等单子。它的代数范畴仍为集合范畴，不能恢复所有拓扑。具体地，同一个二点集从离散拓扑到平凡拓扑的恒等映射是连续双射，却不是同胚；下节的判据会指出这个例子在哪项假设上失败。
\symidx{\(\mathbf{Top}\)：拓扑空间及连续映射组成的范畴}

% !TeX root = ../../Book4.tex
% 纲要标题：### A.3 Barr–Beck 定理的假设与典型例子
% 纲要位置：工作记录/计划纲要.md，第 3836 行。
% 纲要细目（写作范围）：
% #### A.3.1 分裂余等化子
% #### A.3.2 Barr–Beck 单子性判据
% #### A.3.3 模的限制标量与单子性
% #### A.3.4 忠实平坦下降的含义


\section{Barr–Beck 定理的假设与典型例子}
\label{b4:sec:A03}

比较函子总能构造，但它何时为等价，需要一个可以检验的条件。基本观察是：每个单子代数都能由两个自由代数之间的余等化子给出，而且遗忘结构后，这个余等化子具有显式分裂。下面只要求原范畴正确处理这种特殊的余等化子。

\subsection{分裂余等化子}
\label{b4:subsec:A03:01}

\begin{definition}\label{b4:def:A-split-coequalizer}
在范畴 \(\mathcal C\) 中，设
\[
A\mathrel{\substack{\xrightarrow{\ f\ }\\[-0.4ex]\xrightarrow[\ g\ ]{}}}B
\xrightarrow{q}C,\qquad qf=qg.
\]
如果还存在 \(s:C\to B\)、\(t:B\to A\)，满足
\begin{equation}\label{b4:eq:A-split-coequalizer}
qs=1_C,\qquad ft=1_B,\qquad gt=sq,
\end{equation}
则称这个叉形图连同 \(s,t\) 为一个\term[fenlie-yudenghuazi]{分裂余等化子}[split coequalizer]。
\end{definition}

等式 \(qs=1_C\)、\(ft=1_B\) 分别给出 \(q,f\) 的截面；还必须满足 \(gt=sq\)，使沿 \(t\) 返回 \(A\) 后再施行 \(g\)，与先施行 \(q\) 再沿 \(s\) 返回相同。若两张平行箭头只有共同截面 \(r:B\rightarrow A\)，已知的只是 \(fr=gr=1_B\)。这与 \(gt=sq\) 所要求的相容性不同，不能用共同截面代替定义中的全部等式。

\begin{lemma}\label{b4:lem:A-split-absolute}
设 \(\mathcal C\) 为范畴，\(f,g:A\rightrightarrows B\)、\(q:B\to C\) 为其中的态射。若 \(qf=qg\)，并存在 \(s:C\to B\)、\(t:B\to A\) 满足 \(qs=1_C\)、\(ft=1_B\)、\(gt=sq\)，则 \(q\) 确实是 \(f,g\) 的余等化子。任意函子 \(H:\mathcal C\to\mathcal E\)（\(\mathcal E\) 为任意范畴）都保持这个余等化子及其分裂。
\end{lemma}
\begin{proof}
若 \(h:B\to Z\) 满足 \(hf=hg\)，则
\[
h=hft=hgt=hsq.
\]
所以 \(hs:C\to Z\) 给出所需分解。若 \(uq=h\)，则 \(u=uqs=hs\)，因而唯一。任意函子都保持复合、恒等及所列等式，对像图应用同一证明即可。
\end{proof}

对函子 \(U:\mathcal D\to\mathcal C\)，如果平行态射 \(f,g:A\rightrightarrows B\) 的像在 \(\mathcal C\) 中可以补成一个分裂余等化子，就称它们是 \(U\)-分裂的。分裂态射只要求存在于 \(\mathcal C\)，不要求来自 \(\mathcal D\)。

\begin{proposition}\label{b4:prop:A-em-split-coequalizers}
设 \((T,\eta,\mu)\) 为局部小范畴 \(\mathcal C\) 上的单子，\(\mathcal C^T\) 为其代数范畴，\(U^T:\mathcal C^T\to\mathcal C\) 为遗忘函子。若 \(\mathcal C^T\) 中一对平行代数同态
\[
f,g:(A,a)\rightrightarrows(B,b)
\]
的底层图在 \(\mathcal C\) 中有分裂余等化子 \(q:B\to C\)，则 \(C\) 上存在唯一代数结构 \(c:TC\to C\)，使 \(q\) 为代数同态；带此结构的 \(q\) 是 \(\mathcal C^T\) 中的余等化子。此外，\(U^T\) 反映同构。
\end{proposition}
\begin{proof}
由分裂等式 \(qs=1_C\)，施行 \(T\) 及 \(T^2\) 得到
\[
Tq\,Ts=1_{TC},\qquad T^2q\,T^2s=1_{T^2C}.
\]
所以 \(q,Tq,T^2q\) 都是分裂满态射；这里不要求 \(T\) 保持一般满态射。\cref{b4:lem:A-split-absolute}又说明 \(Tq\) 是 \(Tf,Tg\) 的余等化子。于是
\[
qbT(f)=qfa=qga=qbT(g),
\]
所以存在唯一 \(c:TC\to C\)，使 \(cTq=qb\)。这正是 \(q\) 成为代数同态的条件。

由于 \(q\) 为满态射，验证 \(c\eta_C=1_C\) 可以在右边复合 \(q\)；此时由单位自然性得到
\[
c\eta_Cq=cTq\eta_B=qb\eta_B=q.
\]
结合律则在右边复合 \(T^2q\)。这个态射也是分裂满态射，且
\[
cT(c)T^2q=qbT(b),\qquad
c\mu_CT^2q=qb\mu_B.
\]
两者由 \(b\) 的代数公理相等，故 \(cT(c)=c\mu_C\)。

若代数同态 \(h:(B,b)\to(D,d)\) 满足 \(hf=hg\)，则底层余等化子的泛性质给出唯一 \(v:C\to D\)，使 \(vq=h\)。复合满态射 \(Tq\) 后，
\[
vcTq=vqb=hb=dT(h)=dT(v)Tq,
\]
所以 \(vc=dT(v)\)，即 \(v\) 仍是代数同态。这证明了余等化子的泛性质。

最后，若代数同态 \(u:(X,a)\to(Y,b)\) 的底层态射可逆，从 \(ua=bT(u)\) 两边乘以逆，得到 \(u^{-1}b=aT(u^{-1})\)。因此底层逆也是代数同态，\(U^T\) 反映同构。
\end{proof}

\subsection{Barr–Beck 单子性判据}
\label{b4:subsec:A03:02}

如果函子把某个态射送成同构就能推出该态射本身为同构，称这个函子\term[baoshou-hanzi]{保守}[conservative]，也称它反映同构。这个条件不同于忠实：忠实只要求每个同态集合上的映射为单射。

\begin{theorem}[Barr–Beck 单子性判据]\label{b4:thm:A-beck-monadicity}
设 \(\mathcal C,\mathcal D\) 为局部小范畴，\(U:\mathcal D\to\mathcal C\) 有左伴随 \(F:\mathcal C\to\mathcal D\)，伴随的单位、余单位为 \(\eta,\varepsilon\)，所生单子为 \(T=UF\)。则比较函子 \(K:\mathcal D\to\mathcal C^T\) 全忠实且本质满，当且仅当下列条件均成立：
\begin{enumerate}
\item \(U\) 反映同构；
\item 每一对 \(U\)-分裂平行态射在 \(\mathcal D\) 中都有余等化子，且 \(U\) 保持该余等化子。
\end{enumerate}
若这些条件成立，并且已为每个 \(T\)-代数 \((X,a)\) 同时指定平行对
\(F(a),\varepsilon_{FX}:FTX\rightrightarrows FX\) 的余等化子
\(q_{(X,a)}:FX\to L_{(X,a)}\)，则这些已选数据给出 \(K\) 的拟逆，因而 \(K\) 是范畴等价。反过来，若已给出 \(K\) 的拟逆，这族余等化子也可由该拟逆同时构造。
\end{theorem}
这是 \term[barr-beck-danzixing]{Barr–Beck 单子性定理}[Barr--Beck monadicity theorem]的分裂余等化子版本\footnote{巴尔（Michael Barr，1937-），美国数学家，研究范畴论与同调代数；贝克（Jonathan Mock Beck，1935-2006），美国数学家，建立以其命名的单子性判据。}。第二项只涉及指定的一类平行态射，而不是要求原范畴存在并保持全部余等化子。
\begin{proof}
先设 \(K\) 全忠实且本质满。由\cref{b4:prop:A-em-split-coequalizers}，\(U^T\) 反映同构；全忠实的 \(K\) 也反映同构，而 \(U=U^TK\)，故 \(U\) 反映同构。

给定 \(\mathcal D\) 中的 \(U\)-分裂平行对 \(f,g:A\rightrightarrows B\)，并在 \(\mathcal C\) 中取其分裂余等化子 \(q:UB\to C\)。同一命题在 \(C\) 上给出代数结构 \(c\)，并使
\(\widehat q:K(B)\to(C,c)\) 成为 \(K(f),K(g)\) 的余等化子，其底层态射就是 \(q\)。本质满性给出对象 \(L\in\mathcal D\) 及同构
\(\psi:K(L)\xrightarrow{\sim}(C,c)\)；全忠实性给出唯一 \(h:B\to L\)，满足
\(K(h)=\psi^{-1}\widehat q\)。它余等化 \(f,g\)。对任何余等化它们的 \(r:B\to Z\)，\(K(r)\) 经 \(\widehat q\) 唯一分解；再复合 \(\psi\)，利用 \(K\) 的全忠实性，得到 \(r\) 经 \(h\) 的唯一分解。因此 \(h\) 是 \(\mathcal D\) 中的余等化子。又
\(q=U^T(\psi)U(h)\)，且 \(U^T(\psi)\) 为同构，故 \(U(h)\) 也是底层余等化子。其他余等化子与 \(h\) 唯一同构，所以也被 \(U\) 保持。这证明第二项；特别地，当 \(K\) 已是等价时，这一构造正是沿等价运输代数范畴中由 \(U^T\) 创建的分裂余等化子。

反过来，假定两项条件成立。先证明每个 \(T\)-代数都由比较函子得到。对 \((X,a)\)，代数结合律比较的是 \(T^2X\) 上的两种求值次序：\(\mu_X\) 先归并两层形式表达式，\(T(a)\) 先对内层求值，最后都再施行 \(a\)。在有限表单子的例子中，它们把 \([[x_1,x_2],[x_3]]\) 分别送到
\[
[x_1,x_2,x_3],\qquad
[a([x_1,x_2]),a([x_3])],
\]
而 \(a\mu_X=aT(a)\) 要求这两张表的最终值相同。因此在自由对象上应当识别的，正是下面两支平行箭头：
\begin{equation}\label{b4:eq:A-free-presentation}
FTX
\mathrel{\substack{\xrightarrow{\ \varepsilon_{FX}\ }\\[-0.4ex]
\xrightarrow[\ F(a)\ ]{}}}
FX.
\end{equation}
应用 \(U\) 后，它们是 \(\mu_X,T(a):T^2X\rightrightarrows TX\)。代数公理说明 \(a:TX\to X\) 余等化这对态射。取 \(s=\eta_X\)、\(t=\eta_{TX}\)，便有
\[
a\eta_X=1_X,\qquad
\mu_X\eta_{TX}=1_{TX},\qquad
T(a)\eta_{TX}=\eta_Xa.
\]
因此这是一个分裂余等化子。由第二项假设，\eqref{b4:eq:A-free-presentation}存在余等化子 \(q:FX\to L\)，且 \(Uq:TX\to UL\) 也是底层余等化子。余等化子的唯一性给出唯一同构
\(\phi:UL\xrightarrow{\sim}X\)，满足 \(\phi Uq=a\)。记
\(c_L=U(\varepsilon_L):T(UL)\rightarrow UL\)。此时只恢复了正确的底层对象，还须确认 \(K(L)\) 的求值结构也与 \(a\) 相同。把其结构经 \(\phi\) 运输到 \(X\)，得到
\[
b=\phi c_LT(\phi^{-1}):TX\longrightarrow X.
\]
余单位关于 \(q\) 的自然性为
\(q\varepsilon_{FX}=\varepsilon_LF(Uq)\)。应用 \(U\) 并在左侧复合 \(\phi\)，得到两个从 \(T^2X\) 到 \(X\) 的态射相等：
\[
 a\mu_X=\phi c_LT(Uq)=bT(a).
\]
代数公理又给出 \(a\mu_X=aT(a)\)，故 \(bT(a)=aT(a)\)。由单位律，
\[
T(a)T(\eta_X)=T(a\eta_X)=1_{TX}.
\]
在两边右复合 \(T(\eta_X)\)，便得 \(b=a\)。于是
\(\phi c_L=aT(\phi)\)，所以 \(\phi\) 是代数同构
\(K(L)\xrightarrow{\sim}(X,a)\)，证明本质满性。

要从底层代数同态恢复 \(\mathcal D\) 中的态射，需要先证明：每个 \(Y\in\mathcal D\) 的余单位 \(\varepsilon_Y\) 就是其自由展示的余等化子。写 \(a_Y=U(\varepsilon_Y)\)。相同构造给出
\[
FTUY\rightrightarrows FUY\xrightarrow{q}Y',
\]
其中两箭头为 \(\varepsilon_{FUY}\)、\(F(a_Y)\)，而 \(Uq\) 是余等化子。余单位自然性说明 \(\varepsilon_Y:FUY\to Y\) 也余等化这对态射，故唯一分解为 \(hq\)。\(U(\varepsilon_Y)=a_Y\) 本身是上面已经验证的分裂余等化子，所以 \(Uh\) 为两个余等化子之间的同构。由第一项假设，\(h\) 为同构。因此 \(\varepsilon_Y\) 就是这一对态射的余等化子，特别是满态射。

若 \(r,s:Y\to Z\) 满足 \(Ur=Us\)，则余单位自然性给出
\[
r\varepsilon_Y=\varepsilon_ZF(Ur)
=\varepsilon_ZF(Us)=s\varepsilon_Y.
\]
消去满态射 \(\varepsilon_Y\)，得到 \(r=s\)，故 \(U\) 以及 \(K\) 忠实。

现在给定代数同态 \(\alpha:K(Y)\to K(Z)\)，即
\(\alpha a_Y=a_ZT(\alpha)\)。令 \(\widehat\alpha=\varepsilon_ZF(\alpha):FUY\to Z\)。两次使用余单位自然性可得
\[
\begin{aligned}
\widehat\alpha\,\varepsilon_{FUY}
&=\varepsilon_Z\varepsilon_{FUZ}FT(\alpha)\\
&=\varepsilon_ZF(a_Z)FT(\alpha)\\
&=\varepsilon_ZF(\alpha a_Y)
=\widehat\alpha\,F(a_Y).
\end{aligned}
\]
因此它经过余等化子 \(\varepsilon_Y\) 唯一分解为 \(r\varepsilon_Y\)，其中 \(r:Y\to Z\)。应用 \(U\)，得到
\[
(Ur)a_Y=a_ZT(\alpha)=\alpha a_Y.
\]
\(a_Y\) 有右逆 \(\eta_{UY}\)，消去后便有 \(Ur=\alpha\)。这证明全性。

若自由展示的余等化子 \(q_{(X,a)}\) 已同时指定，刚才的运输给出唯一的代数同构
\(K(L_{(X,a)})\xrightarrow{\sim}(X,a)\)。\cref{b4:thm:equivalence-criterion}中带有已选代表的判据便给出拟逆；它在态射上由 \(K\) 的全忠实性唯一确定。对于大范畴，所需代表由规定的余等化子构造或上述已选数据给出。

反过来，若已有拟逆 \(J:\mathcal C^T\to\mathcal D\) 及自然同构
\(\alpha:1_{\mathcal D}\xrightarrow{\sim}JK\)，对代数 \((X,a)\) 取
\[
q_{(X,a)}=J(a)\alpha_{FX}:FX\longrightarrow J(X,a).
\]
这里 \(a:F^T(X)\to(X,a)\) 是自由代数展示的余等化子，而 \(KF=F^T\)。拟逆保持余等化子，\(\alpha\) 的自然性又识别两条平行箭头，故这些公式同时给出所需余等化子。
\end{proof}

离散拓扑伴随的反例现在有了明确解释：底集合函子不反映同构，连续双射可能不是同胚。对其他例子，即使已核验保守性，也仍须检查第二项指定的余等化子条件。

\subsection{模的限制标量与单子性}
\label{b4:subsec:A03:03}

\begin{proposition}\label{b4:prop:A-restriction-scalars-monadic}
设 \(R\to S\) 为含幺环同态，以下模范畴均以幺性左模及模同态为对象和态射。限制标量函子
\[
U:S\text{-}\mathbf{Mod}\longrightarrow R\text{-}\mathbf{Mod}
\]
具有单子性。相应单子在幺性左 \(R\) 模 \(M\) 上为 \(T(M)=S\otimes_RM\)，其中 \(S\) 的右 \(R\) 作用与所得张量积的左 \(R\) 作用均由 \(R\to S\) 给出；单位为 \(m\mapsto1\otimes m\)，乘法为 \(s\otimes t\otimes m\mapsto st\otimes m\)。
\end{proposition}
\begin{proof}
左伴随为扩张标量 \(F(M)=S\otimes_RM\)：映射
\[
\operatorname{Hom}_S(S\otimes_RM,N)\longrightarrow
\operatorname{Hom}_R(M,UN),\qquad
f\longmapsto\bigl(m\mapsto f(1\otimes m)\bigr)
\]
的逆为 \(h\mapsto\bigl(s\otimes m\mapsto s\cdot h(m)\bigr)\)。平衡关系及 \(R\)-线性保证逆映射良定义。

一个 \(S\)-线性映射若在底层为可逆 \(R\)-线性映射，其逆仍保持 \(S\)-作用，所以 \(U\) 反映同构。对 \(S\)-线性平行态射 \(f,g:M\rightrightarrows N\)，余等化子为
\[
N/\operatorname{im}(f-g).
\]
像本来就是 \(S\)-子模；遗忘到 \(R\) 后，底层像与商均不变，仍是 \(R\)-模中的余等化子。因此全部余等化子都存在并被 \(U\) 保持。上述规定的商模构造同时给出了自由展示的余等化子，故\cref{b4:thm:A-beck-monadicity}给出范畴等价。
\end{proof}

这里的代数结构也能直接辨认。对 \(R\)-模 \(M\)，一个 \(T\)-代数态射 \(a:S\otimes_RM\to M\) 定义 \(s\cdot m=a(s\otimes m)\)。单位律给出 \(1m=m\)，结合律给出 \((st)m=s(tm)\)，平衡关系使它延伸原来的 \(R\)-作用。反向由任何 \(S\)-模作用都得到这样的 \(a\)。这个直接说明与比较函子恰好一致。

\subsection{忠实平坦下降的含义}
\label{b4:subsec:A03:04}

限制标量的单子性对任意环同态成立。下降研究另一个方向：从已经扩张标量的对象以及相容数据，恢复原来的对象。因此这里使用扩张标量函子的余单子，而不是把上一命题再改写一次。

\begin{definition}\label{b4:def:A-faithfully-flat-descent-data}
设 \(R\to S\) 为交换环同态，\(M\) 为 \(S\)-模。若 \(S\)-线性映射
\(\gamma:M\to S\otimes_RM\) 满足
\begin{equation}\label{b4:eq:A-module-descent-data}
\begin{aligned}
e\gamma&=1_M,\qquad e(s\otimes m)=sm,\\
(1_S\otimes\gamma)\gamma&=\delta\gamma,\qquad
\delta(s\otimes m)=s\otimes1\otimes m,
\end{aligned}
\end{equation}
则称 \((M,\gamma)\) 为一个以余单子形式给出的\term[mo-xiajiang-shuju]{模下降数据}[module descent datum]。这里 \(S\) 在张量积的第一个因子上作用。对两份数据 \((M,\gamma)\)、\((M',\gamma')\)，若 \(S\) 线性映射 \(f:M\to M'\) 满足 \((1_S\otimes f)\gamma=\gamma'f\)，则称 \(f\) 为它们之间的同态。
\end{definition}

若 \(S=R\times R\)，一个 \(S\)-模分成两块 \(M_1\oplus M_2\)，而从 \(R\)-模 \(V\) 扩张得到的是 \(V\oplus V\)。因此要判断给定的两块是否来自同一个模，还须指定它们怎样互相识别。\cref{b4:prop:A-product-descent-model}将证明，\(\gamma\) 恰好记录一个同构 \(M_1\xrightarrow{\sim}M_2\)；余结合律使两方向的识别互逆。

\ProseParagraphBreak

这些条件正是伴随 \(S\otimes_R-\dashv U\) 所生余单子的余代数公理。对 \(R\)-模 \(V\)，扩张模 \(S\otimes_RV\) 总带有这样的数据：
\[
\gamma_V(s\otimes v)=s\otimes(1\otimes v).
\]
等式逐个在纯张量上成立。问题是给定的数据是否都这样产生，以及保持数据的态射是否也来自唯一的 \(R\)-线性映射。

\begin{theorem}[忠实平坦模下降]\label{b4:thm:A-faithfully-flat-descent}
设 \(R\to S\) 为交换含幺环同态，且 \(S\) 是忠实平坦 \(R\) 模，即 \(S\otimes_R-\) 正合，并对每个 \(R\) 模 \(V\) 满足
\(S\otimes_RV=0\Rightarrow V=0\)。以 \(\mathsf{Desc}(S/R)\) 表示 \(S/R\) 模下降数据及其同态组成的范畴。令 \(\gamma_V(s\otimes v)=s\otimes(1\otimes v)\)，则
\[
V\longmapsto(S\otimes_RV,\gamma_V)
\]
是 \(R\) 模范畴到 \(\mathsf{Desc}(S/R)\) 的等价，态射 \(f:V\to V'\) 被送到 \(1_S\otimes f\)。给定下降数据 \((M,\gamma)\)，一个恢复对象为在 \(R\)-模范畴中取的等化子
\begin{equation}\label{b4:eq:A-descent-equalizer}
V=\bigl\{m\in M:\gamma(m)=1\otimes m\bigr\},
\end{equation}
相应同构 \(S\otimes_RV\to M\) 为 \(s\otimes v\mapsto sv\)。
\end{theorem}
\symidx{\(\mathsf{Desc}(S/R)\)：模下降数据组成的范畴}
\begin{proof}
把\cref{b4:thm:A-beck-monadicity}应用于对偶范畴：左伴随 \(F=S\otimes_R-\) 具有余单子性，只须它反映同构，并保持所有在 \(F\) 下分裂的等化子。平坦性使 \(F\) 正合，从而保持模范畴中的全部等化子。若 \(F(f)\) 可逆，则正合性给出
\[
F(\ker f)=0,\qquad F(\operatorname{coker}f)=0.
\]
忠实性假设使原来的核、余核都为零，故 \(f\) 可逆。模范畴的等化子可统一取为两态射之差的核子模，给出对偶判据所需的已选对象。因此该判据得到范畴等价。

为说明所列恢复公式，记 \(\eta_M(m)=1\otimes m\)。\(\gamma\) 与 \(\eta_M\) 都是 \(R\)-线性的，但 \(\eta_M\) 一般不是 \(S\)-线性的：
\[
\eta_M(sm)=1\otimes sm,\qquad
s\cdot\eta_M(m)=s\otimes m,
\]
这两个张量未必相同。因此\eqref{b4:eq:A-descent-equalizer}首先是 \(R\)-模中的等化子，给恢复对象 \(V\) 以 \(R\)-模结构。平坦性使扩张后的 \(S\otimes_RV\) 成为 \(S\)-模中的等化子：
\[
1_S\otimes\gamma,\ 1_S\otimes\eta_M:
S\otimes_RM\rightrightarrows S\otimes_RS\otimes_RM
\]
两支箭头都保持第一个 \(S\) 因子上的作用。另一方面，\eqref{b4:eq:A-module-descent-data}说明 \(\gamma\) 的像属于这个等化子，且 \(e\gamma=1_M\)。

若 \(z\in S\otimes_RM\) 满足两映射取值相同，对该等式施行
\(s\otimes t\otimes m\mapsto st\otimes m\)。由于 \(\gamma\) 是 \(S\)-线性的，左边变成 \(\gamma\bigl(e(z)\bigr)\)，右边变成 \(z\)。因此 \(z=\gamma\bigl(e(z)\bigr)\)，说明 \(\gamma\) 把 \(M\) 同构到该等化子。将它与 \(S\otimes_RV\) 的等化子识别相比较，逆映射为
\(\phi:S\otimes_RV\to M\)、\(\phi(s\otimes v)=sv\)。它还保持下降数据：对 \(v\in V\)，
\[
\gamma\phi(s\otimes v)=\gamma(sv)=s\otimes v
=(1_S\otimes\phi)\gamma_V(s\otimes v).
\]
因此 \(\phi\) 是下降数据之间的同构。

对保持下降数据的态射 \(f:(M,\gamma)\to(M',\gamma')\)，若 \(v\in V\)，则
\(\gamma'f(v)=(1_S\otimes f)\gamma(v)=1\otimes f(v)\)，故 \(f\) 限制为恢复子模之间的 \(R\)-线性映射 \(f|_V:V\to V'\)。限制保持恒等与复合，而且
\(\phi'(1_S\otimes f|_V)=f\phi\)，所以所给恢复公式及其同构对态射自然。

最后，对原来的 \(R\)-模 \(W\)，令 \(V_W\) 为标准下降数据 \((S\otimes_RW,\gamma_W)\) 的恢复子模。映射 \(\theta_W:W\rightarrow V_W\)、\(w\mapsto1\otimes w\) 扩张标量后，与刚才的乘法同构复合为 \(1_{S\otimes_RW}\)，故 \(S\otimes_R\theta_W\) 可逆。前面的等化子识别用平坦性得到了 \(S\otimes_RV_W\cong S\otimes_RW\)；此处还要用忠实性，才能从扩张后核、余核消失推出 \(\theta_W\) 本身可逆。这给出另一方向的自然同构，因而上述恢复函子就是所列等价的拟逆。
\end{proof}

对角映射 \(R\rightarrow R\times R\) 是忠实平坦的，因为张量函子把 \(V\) 送到 \(V\oplus V\)，既保持正合列，也不会把非零模变为零。下面把这个例子的下降条件直接写成两块之间的映射。

\begin{proposition}\label{b4:prop:A-product-descent-model}
设 \(R\) 为交换含幺环，\(S=R\times R\)，\(R\rightarrow S\) 为对角映射。令 \(e_1=(1,0)\)、\(e_2=(0,1)\)。对幺性 \(S\)-模 \(M\)，写 \(M=M_1\oplus M_2\)，其中 \(M_i=e_iM\)。在同构 \(S\otimes_RM\cong M\oplus M\) 下，\(M\) 上的模下降数据与 \(R\)-模同构 \(f:M_1\rightarrow M_2\) 一一对应，相应结构为
\[
\gamma(m_1,m_2)
=\Bigl(\bigl(m_1,f(m_1)\bigr),
\bigl(f^{-1}(m_2),m_2\bigr)\Bigr).
\]
其恢复模为 \(f\) 的图
\[
V=\{\,(m_1,f(m_1))\mid m_1\in M_1\,\},
\]
以坐标投影同构于 \(M_1\)。
\end{proposition}
\begin{proof}
同构 \(S\otimes_RM\cong M\oplus M\) 把 \(e_1\otimes m\)、\(e_2\otimes m\) 分别送到 \((m,0)\)、\((0,m)\)；外层两块承受 \(S\) 的两个坐标作用。\(\gamma\) 的 \(S\)-线性因此给出 \(R\)-线性映射 \(\gamma_{ij}:M_j\rightarrow M_i\)，满足
\[
\gamma(m_1,m_2)
=\bigl((\gamma_{11}(m_1),\gamma_{21}(m_1)),
(\gamma_{12}(m_2),\gamma_{22}(m_2))\bigr).
\]
余单位 \(e\) 从外层第一块取内部 \(M_1\) 分量，从第二块取 \(M_2\) 分量，所以 \(e\gamma=1_M\) 等价于 \(\gamma_{11}=1_{M_1}\)、\(\gamma_{22}=1_{M_2}\)。

对 \(m\in M_j\)，比较余结合等式中外层 \(e_j\)、中间层 \(e_i\) 及内部 \(M_k\) 的分量：左边先用 \(\gamma_{ij}\)，再用 \(\gamma_{ki}\)；右边由 \(\delta(e_j\otimes m)=e_j\otimes1\otimes m\) 把中间层同时展开为 \(e_1+e_2\)。因此余结合律恰好要求
\[
\gamma_{ki}\gamma_{ij}=\gamma_{kj}
\qquad(i,j,k\in\{1,2\}).
\]
在两个对角映射已为恒等的条件下，唯一尚需验证的等式是
\[
\gamma_{12}\gamma_{21}=1_{M_1},\qquad
\gamma_{21}\gamma_{12}=1_{M_2}.
\]
故 \(f=\gamma_{21}\) 必为同构，且 \(\gamma_{12}=f^{-1}\)。反过来，任取这样的同构，所列公式是 \(S\)-线性的，并满足刚才全部分量等式，因而给出下降数据。

最后，\(1\otimes(m_1,m_2)\) 在外层两块中都是 \((m_1,m_2)\)。恢复条件 \(\gamma(m_1,m_2)=1\otimes(m_1,m_2)\) 因此等价于 \(m_2=f(m_1)\)，给出所述图子模。
\end{proof}

平坦而不忠实的映射则可能保留相容方程，却丢失原对象。例如 \(\mathbb Z\rightarrow\mathbb Q\) 把所有有限交换群张量成零；不同的原模于是具有相同的零下降数据，无法由它们唯一恢复。

% !TeX root = ../../Book4.tex
% 纲要标题：### A.4 可表性与伴随函子定理
% 纲要位置：工作记录/计划纲要.md，第 3557 行。
% 纲要细目（写作范围）：
% #### A.4.1 可表性与泛元素
% #### A.4.2 解集条件与伴随函子定理
% ---


\section{可表性与伴随函子定理}
\label{b4:sec:A04}

前面的自由代数从一个已经给定的单子构造出来。更一般地，面对一个遗忘函子，是否存在自由对象本身可能就是问题。本节把这个存在性问题化为带有指定元素的对象是否有始对象，并说明小性条件在其中的作用。

\subsection{可表性与泛元素}
\label{b4:subsec:A04:01}

设 \(\mathcal C\) 为局部小范畴，\(H:\mathcal C\to\mathbf{Set}\) 为协变函子。它的\term[yuansu-fanchou]{元素范畴}[category of elements]在这里采用如下方向：对象为 \((X,x)\)，其中 \(x\in H(X)\)；态射
\((X,x)\to(Y,y)\) 为满足 \(H(f)(x)=y\) 的 \(f:X\to Y\)。复合由 \(H\) 的函子性保证，记此范畴为 \(\int H\)。

\begin{proposition}\label{b4:prop:A-representability-initial}
设 \(\mathcal C\) 为局部小范畴，\(H:\mathcal C\to\mathbf{Set}\) 为函子。记 \(\int H\) 为以 \((X,x)\)（\(X\in\mathcal C,x\in H(X)\)）为对象、以满足 \(H(f)(x)=y\) 的态射 \(f:X\to Y\) 为 \((X,x)\to(Y,y)\) 的态射的元素范畴。则 \(H\) 可表，当且仅当 \(\int H\) 有始对象。具体地，若 \((A,u)\) 为始对象，则
\[
\operatorname{Hom}_{\mathcal C}(A,X)\longrightarrow H(X),
\qquad f\longmapsto H(f)(u)
\]
为自然双射。
\end{proposition}
\begin{proof}
从 \((A,u)\) 到 \((X,x)\) 存在唯一态射，恰好意味着每个 \(x\in H(X)\) 唯一写成 \(H(f)(u)\)。这就是\cref{b4:prop:universal-element}的泛元素判据。其自然性也直接来自
\(H(g)H(f)(u)=H(gf)(u)\)。反向由一个表示取 \(u\) 为恒等态射对应的元素，就得到同样的存在唯一性。
\end{proof}

例如固定交换环 \(R\) 及 \(a\in R\)，在交换 \(R\)-代数范畴上研究
\[
H(A)=\{x\in A:x^2=a\,1_A\}.
\]
环 \(R[t]/(t^2-a)\) 连同 \(\bar t\) 给出泛元素：把 \(\bar t\) 送到 \(x\) 的 \(R\)-代数同态存在，当且仅当 \(x^2=a\,1_A\)，而多项式代入使它唯一。这一可表性并不声称每个 \(A\) 中都有平方根；没有根时，相应同态集合也为空。

\ProseParagraphBreak

由\cref{b4:prop:representables-preserve-limits}，可表函子保持所有已经存在的小极限。因此它必须把终对象送到单元素集合，并保持积与等化子。这里的条件是函子条件，而不是逐个比较集合基数。例如取直接像的幂集函子把单元素集合送到一个二元素集合，所以它不可协变表示；取逆像的反变幂集函子却由二元素集合表示。改变变性会改变整个问题。

\subsection{解集条件与伴随函子定理}
\label{b4:subsec:A04:02}

一般范畴的全部对象可能组成真类，不能未经检查就把它们全部取积。伴随函子定理中的解集条件，正是要把足以产生全部解的对象控制在一个集合内。

\begin{lemma}\label{b4:lem:A-weakly-initial-set}
设 \(\mathcal E\) 为局部小、具有全部小极限的范畴。如果存在由一个集合指标化的对象族 \((E_i)_{i\in I}\)，使每个 \(X\in\mathcal E\) 至少接收一个来自某个 \(E_i\) 的态射，那么 \(\mathcal E\) 有始对象。
\end{lemma}
这样的族称为一个\term[ruoshi-duixiangji]{弱始对象集}[weakly initial set]；“弱”表示只要求态射存在，尚不要求唯一。
\begin{proof}
取积 \(P=\prod_{i\in I}E_i\)。对任意 \(X\)，先选一个 \(E_i\to X\)，再与积投影 \(P\to E_i\) 复合，得到 \(P\to X\)。所以 \(P\) 是弱始对象。

局部小性保证 \(\operatorname{End}_{\mathcal E}(P)\) 是一个集合。取全部对 \((f,1_P)\) 的共同等化子 \(e:E\to P\)，即
\[
fe=e\qquad\bigl(f\in\operatorname{End}_{\mathcal E}(P)\bigr).
\]
这可由集合指标的积与一个等化子构造，故存在。任意 \(X\) 都接收 \(E\to P\to X\)，所以只须证明唯一性。

给定 \(u,v:E\rightrightarrows X\)，取等化子 \(q:Y\to E\)。由 \(P\) 弱始，有 \(r:P\to Y\)。复合 \(eqr\) 是 \(P\) 的自同态，所以
\[
eqre=e.
\]
消去单态射 \(e\)，得到 \(qre=1_E\)。因此 \(q\) 同时为单态射和分裂满态射，是同构。由 \(uq=vq\) 得 \(u=v\)。这就证明 \(E\) 为始对象。
\end{proof}

\begin{definition}\label{b4:def:A-solution-set}
设 \(\mathcal C,\mathcal D\) 为局部小范畴，\(U:\mathcal D\to\mathcal C\) 为函子。若对每个 \(X\in\mathcal C\)，都存在由集合 \(I_X\) 指标的对象族 \((Y_i)_{i\in I_X}\)（\(Y_i\in\mathcal D\)）与态射族
\[
\{u_i:X\to UY_i\}_{i\in I_X},
\]
使对每个 \(Y\in\mathcal D\) 及态射 \(u:X\to UY\)，都存在 \(i\in I_X\) 与态射 \(v:Y_i\to Y\)，满足 \(u=U(v)u_i\)，则称 \(U\) 满足\term[jieji-tiaojian]{解集条件}[solution set condition]。
\end{definition}

分解不要求唯一，所选族也可以随 \(X\) 改变。要求的是存在一个集合大小的候选族，足以经由后续态射产生所有解。

\begin{theorem}[一般伴随函子定理]\label{b4:thm:A-general-adjoint-functor}
设 \(\mathcal C,\mathcal D\) 为局部小范畴，\(\mathcal D\) 有全部小极限，\(U:\mathcal D\to\mathcal C\) 为函子。假定当 \(U\) 满足解集条件时，所需的各解集及小极限可以按对象类同时选择；可把这些选择作为已给数据，或采用足以保证这些选择的类选择公理。在这一选择前提下，\(U\) 有左伴随，当且仅当它保持全部小极限并满足解集条件。
\end{theorem}
这一形式称为 \term[freyd-bansui-hanzi-dingli]{Freyd 一般伴随函子定理}[Freyd's general adjoint functor theorem]。\footnote{弗雷德（Peter John Freyd，1936-），美国数学家，研究范畴论，提出伴随函子定理与 Freyd--Mitchell 嵌入定理。}
\begin{proof}
若 \(F\dashv U\)，\cref{b4:thm:adjoints-preserve-limits}说明 \(U\) 保持小极限。对每个 \(X\)，伴随单位这一个态射 \(\eta_X:X\to UFX\) 已经满足解集条件：每个 \(X\to UY\) 唯一经它分解。

反过来，固定 \(X\in\mathcal C\)。考虑逗号范畴 \((X\downarrow U)\)：对象为 \((Y,u:X\to UY)\)，态射为使 \(U(v)u=u'\) 的 \(v:Y\to Y'\)。它局部小，因为每个态射集合是 \(\mathcal D\) 中某个态射集合的子集。它也有全部小极限。确实，对一个小图，先在 \(\mathcal D\) 中取底层对象的极限 \(L\)；由于 \(U\) 保持此极限，图中各个 \(X\to UY\) 唯一合成 \(X\to UL\)。极限锥与这个态射便满足逗号范畴的极限泛性质。

解集条件恰好说 \((X\downarrow U)\) 有弱始对象集。使用假定中可同时选定的解集及所需小极限，按\cref{b4:lem:A-weakly-initial-set}的构造，为每个 \(X\in\mathcal C\) 同时指定始对象，记为
\((FX,\eta_X:X\to UFX)\)。其泛性质是：任意 \(u:X\to UY\) 都唯一写成 \(U(f)\eta_X\)，其中 \(f:FX\to Y\)。

对 \(a:X\to X'\)，用这个泛性质唯一规定 \(F(a):FX\to FX'\)，使
\[
U\bigl(F(a)\bigr)\eta_X=\eta_{X'}a.
\]
唯一性说明 \(F(1_X)=1_{FX}\)、\(F(ba)=F(b)F(a)\)，所以 \(F\) 成为函子。上述分解还给出自然双射
\(\operatorname{Hom}_{\mathcal D}(FX,Y)\cong
\operatorname{Hom}_{\mathcal C}(X,UY)\)，因而 \(F\dashv U\)。
\end{proof}

例如群的底集合函子 \(U:\mathbf{Grp}\to\mathbf{Set}\) 保持小极限：积与等化子上的群运算逐坐标定义，底集合正是集合极限；\cref{b4:thm:limits-products-equalizers}再给出全部小极限。对固定集合 \(X\)，令
\(\kappa=\max(|X|,\aleph_0)\)。任意映射 \(X\to UG\) 的像生成一个基数至多为 \(\kappa\) 的子群，因为它的元素由有限长的生成元及其逆的词给出。把所有基数不超过 \(\kappa\) 的群搬到某个固定集合的子集上，其可能的乘法、逆及映射 \(X\to G\) 组成一个集合。这些映射给出解集：原来的映射先进入像生成的子群，再包含到 \(G\)。因此定理证明自由群存在，并同时给出其泛性质；具体的既约词模型是进一步的构造。

\begin{corollary}\label{b4:cor:A-representability-solution-set}
设 \(\mathcal C\) 为局部小、具有全部小极限的范畴，\(H:\mathcal C\to\mathbf{Set}\) 为保持小极限的函子。若存在一个集合
\(\bigl\{(A_i,x_i):x_i\in H(A_i)\bigr\}_{i\in I}\)，使每个 \(x\in H(X)\) 都形如 \(H(f)(x_i)\)，其中 \(f:A_i\to X\)，则 \(H\) 可表。反过来，可表函子满足这些条件。
\end{corollary}
\begin{proof}
保持极限说明元素范畴 \(\int H\) 有全部小极限：先在 \(\mathcal C\) 中取极限，再利用 \(H\) 保持该极限，把相容元素族合成为一个元素。假设中的集合是它的弱始对象集，故\cref{b4:lem:A-weakly-initial-set}给出始对象，再由\cref{b4:prop:A-representability-initial}得到可表性。

反过来，可表函子保持小极限，而表示对象连同其泛元素这一个对象对，已经给出所需的集合。
\end{proof}

因此，保持极限解决相容性问题，解集条件解决大小问题，始对象的唯一性再给出泛性质。使用这类定理时，应分别指出这三步的根据，而不能把“保持极限”单独作为存在自由对象或表示对象的充分理由。

% !TeX root = ../../Book4.tex
% 纲要标题：### A.5 Kleisli 范畴与代数理论
% 纲要细目（写作范围）：
% ---
% #### A.5.1 Kleisli 态射与替换的复合
% #### A.5.2 Kleisli 伴随与自由代数
% #### A.5.3 有限积与代数理论的模型

\section{Kleisli 范畴与代数理论}
\label{b4:sec:A05}

把每个字母替换成一个字，便得到自由幺半群之间的同态。若再把新字母替换成字，复合时必须把内层的字展开、连接，而不是把它们当成不能拆开的新字母。单子的单位与乘法恰好提供这种替换的恒等与复合。与把表达式送到实际值的单子代数相比，这里研究的是表达式到表达式的映射；随后还可以只保留有限多个变量，把运算及其恒等式组织成一个带有限积的范畴。

\subsection{Kleisli 态射与替换的复合}
\label{b4:subsec:A05:01}

\begin{definition}\label{b4:def:A-kleisli-category}
设 \((T,\eta,\mu)\) 为局部小范畴 \(\mathcal C\) 上的单子。令 \(\mathcal C_T\) 的对象与 \(\mathcal C\) 相同，并规定
\[
\operatorname{Hom}_{\mathcal C_T}(X,Y)
=\operatorname{Hom}_{\mathcal C}(X,TY).
\]
对 \(f:X\to TY\)、\(g:Y\to TZ\)，规定复合及恒等为
\[
g\star f=\mu_Z\circ T(g)\circ f,
\qquad 1_X^{\mathcal C_T}=\eta_X.
\]
所得范畴称为 \(T\) 的\term[kleisli-fanchou]{Kleisli 范畴}[Kleisli category]。
\end{definition}

同一个 \(f\) 在 \(\mathcal C_T\) 中的目标是 \(Y\)，在 \(\mathcal C\) 中的目标却是 \(TY\)。复合的中间步骤为
\[
X\xrightarrow{f}TY\xrightarrow{T(g)}T^2Z
\xrightarrow{\mu_Z}TZ.
\]
例如有限表单子中，\(f\) 给每个 \(x\) 一个由 \(Y\) 中元素组成的表，\(T(g)\) 把每项换成一个表，\(\mu_Z\) 再把这些表连接起来。

\begin{proposition}\label{b4:prop:A-kleisli-laws}
设 \((T,\eta,\mu)\) 为局部小范畴 \(\mathcal C\) 上的单子。上述复合满足结合律，两端的 \(\eta\) 为恒等态射，因而 \(\mathcal C_T\) 是局部小范畴。
\end{proposition}
\begin{proof}
给定 \(f:X\to TY\)、\(g:Y\to TZ\)、\(h:Z\to TW\)。乘法关于 \(h\) 的自然性给出
\(T(h)\mu_Z=\mu_{TW}T^2(h)\)，再用单子的结合律，得到
\[
\begin{aligned}
h\star(g\star f)
&=\mu_WT(h)\mu_ZT(g)f\\
&=\mu_W\mu_{TW}T^2(h)T(g)f\\
&=\mu_WT(\mu_W)T^2(h)T(g)f\\
&=\mu_WT\bigl(\mu_WT(h)g\bigr)f
=(h\star g)\star f.
\end{aligned}
\]
对恒等态射，单位律给出
\[
\eta_Y\star f=\mu_YT(\eta_Y)f=f.
\]
单位的自然性 \(T(f)\eta_X=\eta_{TY}f\) 以及另一条单位律则给出
\[
f\star\eta_X=\mu_YT(f)\eta_X
=\mu_Y\eta_{TY}f=f.
\]
每个 Hom 集都是 \(\mathcal C\) 中的一个 Hom 集，因此局部小。
\end{proof}

\subsection{Kleisli 伴随与自由代数}
\label{b4:subsec:A05:02}

普通态射 \(u:X\to Y\) 可以先映入 \(TY\)，于是成为 Kleisli 态射。反方向则把一个替换 \(f:X\to TY\) 延伸到全部形式表达式 \(TX\)：先逐项替换，再用单子乘法归并。

\begin{theorem}\label{b4:thm:A-kleisli-adjunction}
设 \((T,\eta,\mu)\) 为局部小范畴 \(\mathcal C\) 上的单子。函子 \(J_T:\mathcal C\to\mathcal C_T\)、\(V_T:\mathcal C_T\to\mathcal C\) 由以下公式定义：
\[
\begin{aligned}
J_T(X)&=X,&J_T(u)&=\eta_Yu\quad(u:X\to Y),\\
V_T(X)&=TX,&V_T(f)&=\mu_YT(f)\quad(f:X\to TY).
\end{aligned}
\]
有 \(J_T\dashv V_T\)，且该伴随产生的单子为原来的 \((T,\eta,\mu)\)。
\end{theorem}
\begin{proof}
对 \(u:X\to Y\)、\(v:Y\to Z\)，单位的自然性和单位律给出
\[
(\eta_Zv)\star(\eta_Yu)
=\mu_ZT(\eta_Z)T(v)\eta_Yu
=\eta_Zvu.
\]
又 \(J_T(1_X)=\eta_X\)，所以 \(J_T\) 为函子。\(V_T\) 把恒等 \(\eta_X\) 送到 \(1_{TX}\)。对 Kleisli 态射 \(f:X\to TY\)、\(g:Y\to TZ\)，有
\[
\begin{aligned}
V_T(g)V_T(f)
&=\mu_ZT(g)\mu_YT(f)\\
&=\mu_Z\mu_{TZ}T^2(g)T(f)\\
&=\mu_ZT(\mu_Z)T^2(g)T(f)
=V_T(g\star f),
\end{aligned}
\]
故 \(V_T\) 也为函子。

根据 Hom 集的定义，双射
\[
\operatorname{Hom}_{\mathcal C_T}(J_TX,Y)
=\operatorname{Hom}_{\mathcal C}(X,TY)
=\operatorname{Hom}_{\mathcal C}(X,V_TY)
\]
就是恒等映射。它对 \(X\) 自然：在 Kleisli 一侧预复合 \(J_T(u)\)，经单位自然性及单位律化为普通的预复合 \(u\)。它对 \(Y\) 也自然：后复合 Kleisli 态射 \(g\)，恰是后复合普通态射 \(V_T(g)=\mu T(g)\)。因此双射给出伴随。

单位在 \(X\) 处为 \(\eta_X\)。余单位在 Kleisli 对象 \(Y\) 处是 \(J_TV_TY\to Y\) 的态射，它在 \(\mathcal C\) 中由 \(1_{TY}:TY\to TY\) 表示。于是 \(V_TJ_T=T\)，而伴随的乘法 \(V_T(\varepsilon_{J_TX})\) 为 \(\mu_XT(1_{TX})=\mu_X\)，恢复原单子。
\end{proof}

\begin{proposition}\label{b4:prop:A-kleisli-free-algebras}
设 \((T,\eta,\mu)\) 为局部小范畴 \(\mathcal C\) 上的单子，\(\mathcal C^T\) 为 Eilenberg–Moore 范畴。公式
\[
H:\mathcal C_T\longrightarrow\mathcal C^T,
\qquad H(X)=(TX,\mu_X),\quad H(f)=\mu_YT(f)
\]
给出全忠实函子，其对象像为全部自由 \(T\)-代数。在各自由代数上指定一个自由生成对象后，它给出 \(\mathcal C_T\) 与自由代数所成全子范畴之间的等价。
\end{proposition}
\begin{proof}
由\cref{b4:thm:A-em-adjunction}，\(\mu_YT(f)\) 是从 \((TX,\mu_X)\) 到 \((TY,\mu_Y)\) 的代数同态。上一证明已经验证恒等与复合相容，所以 \(H\) 为函子。自由代数的伴随双射为
\[
\operatorname{Hom}_{\mathcal C^T}(H(X),H(Y))
\cong\operatorname{Hom}_{\mathcal C}(X,TY),
\]
右侧正是 Kleisli 态射集，逆双射为 \(f\mapsto\mu_YT(f)\)。因此 \(H\) 全忠实。它的对象正是 \((TX,\mu_X)\)，所以像中包含所有自由代数；指定各自由生成对象后，用上述 Hom 双射恢复态射便得到拟逆。
\end{proof}

对有限表单子，\(\mathbf{Set}_T\) 的态射 \(X\to Y\) 给每个 \(X\) 字母指定一个 \(Y\) 字。上述函子把它延伸成自由幺半群之间的同态。单子代数范畴却包含任意幺半群，而不只包含自由幺半群。因此，Kleisli 范畴描述自由对象之间的替换，Eilenberg–Moore 范畴描述允许求值的全部代数；二者的对象范围不同。相关的单子与伴随构造可对照 Mac Lane\cite[Chapter VI]{MacLane1998Categories}。

\subsection{有限积与代数理论的模型}
\label{b4:subsec:A05:03}

一个有 \(n\) 个变量的群词，可以在任意群中给出一个 \(n\) 元运算。把一组词代入另一组词，就是运算的复合；把各变量分别取出，就是投影。只考虑这些词和代入，仍能保留群运算的全部恒等式，而不必预先指定一个群。Lawvere 的代数理论用有限积把这一结构表达为范畴。

\begin{definition}\label{b4:def:A-lawvere-theory}
设 \(\mathcal L\) 为小范畴，其对象标为非负整数 \(0,1,2,\ldots\)。若指定对象 \(1\)，并为每个 \(n\) 指定把 \(n\) 表示为 \(1\) 的 \(n\) 次积的投影
\[
\pi_i^n:n\longrightarrow1\quad(1\le i\le n),
\]
其中 \(0\) 为终对象，\(\pi_1^1=1_1\)，则称这组数据为单类型的\term[lawvere-daishu-lilun]{Lawvere 代数理论}[Lawvere algebraic theory]。这里“单类型”表示每个模型只有一种底集合。一个保持有限积的函子 \(M:\mathcal L\to\mathbf{Set}\) 称为该理论的模型。模型以自然变换为态射，组成范畴 \(\operatorname{Mod}(\mathcal L)\)。
\end{definition}

各 \(n\) 都是 \(1\) 的有限次积，故 \(n+m\) 以分块投影成为 \(n,m\) 的积，\(\mathcal L\) 确有全部有限积。一个态射 \(a:n\to1\) 表示一个有 \(n\) 个变量的运算项；\(0\to1\) 表示常数项。任意 \(f:n\to m\) 由其分量 \(\pi_j^mf:n\to1\) 唯一决定，因而是一组 \(m\) 个运算项。若 \(g:m\to r\)，则 \(gf\) 的分量是把 \(f\) 的各项代入 \(g\) 的各项。两个项在 \(\mathcal L\) 中相等，就要求它们在每个模型中取相同值；关系由态射的等式表达。

\begin{proposition}\label{b4:prop:A-lawvere-model-operations}
设 \(\mathcal L\) 为单类型 Lawvere 理论，\(M\) 为其集合模型，并令 \(S=M(1)\)。积投影给出典范双射 \(M(n)\cong S^n\)，其中 \(S^0\) 为单元素集合。因此 \(M\) 的数据等价于为每个 \(a:n\to1\) 指定函数 \(a_S:S^n\to S\)，使投影成为坐标投影，并使项的复合成为函数的代入。模型的自然变换恰好对应保持全部这些运算的函数。
\end{proposition}
\begin{proof}
保持积给出典范双射
\[
\kappa_n:M(n)\xrightarrow{\sim}S^n,
\qquad z\longmapsto\bigl(M(\pi_i^n)(z)\bigr)_i.
\]
用它运输 \(M(a)\) 得到 \(a_S=M(a)\kappa_n^{-1}\)。函子公理保证投影与代入等式。

反过来，设给定这样的运算。定义 \(M(n)=S^n\)；对 \(f:n\to m\)，定义
\[
M(f)(x)=\bigl((\pi_j^mf)_S(x)\bigr)_{j=1}^m.
\]
当 \(m=0\) 时取唯一函数。投影公理使 \(M(1_n)\) 为恒等，代入公理使 \(M(gf)=M(g)M(f)\)。分块坐标还保证它保持有限积，故确实给出模型。

若 \(\alpha:M\Rightarrow N\)，对投影的自然性迫使 \(\kappa_n^N\alpha_n(\kappa_n^M)^{-1}=(\alpha_1)^n\)。对每个 \(a:n\to1\) 的自然性遂等价于
\[
\alpha_1\bigl(a_S(x_1,\ldots,x_n)\bigr)
=a_{N(1)}\bigl(\alpha_1(x_1),\ldots,\alpha_1(x_n)\bigr).
\]
反过来，保持全部运算的函数按各次幂定义分量，上述等式使它对每个分量项、从而对任意态射自然。于是对象与态射的两种描述都互相恢复。
\end{proof}

以群为例，令 \(F_n\) 为具有指定自由生成元 \(x_1,\ldots,x_n\) 的自由群，\(F_0\) 为平凡群。取这些自由群及全部群同态组成的范畴的相反范畴 \(\mathcal L_{\mathrm{Grp}}\)，即
\[
\operatorname{Hom}_{\mathcal L_{\mathrm{Grp}}}(n,m)
=\operatorname{Hom}_{\mathbf{Grp}}(F_m,F_n).
\]
一个这样的同态由 \(m\) 个 \(n\) 变量群词决定。自由群的泛性质给出 \(F_{n+m}=F_n*F_m\)，所以在相反范畴中 \(n+m\) 是积；\(F_0\) 是始对象，所以 \(0\) 是终对象。这是一项 Lawvere 理论。

\ProseParagraphBreak
在其模型的底集合 \(S\) 上，词 \(x_1x_2:2\to1\)、\(x_1^{-1}:1\to1\) 与空词 \(0\to1\) 分别定义乘法、逆元与单位。自由群中的词等式给出结合律及单位、逆元等式，故 \(S\) 是群。反过来，任意群 \(G\) 都能按各群词的求值定义模型 \(M_G(n)=G^n\)，代入与求值相容。一般模型中的每个词也由这三个基本运算计算，因而与从所得群构造的模型通过 \(\kappa_n\) 自然同构。保持这三个运算的函数恰为群同态，遂得到
\[
\operatorname{Mod}(\mathcal L_{\mathrm{Grp}})\simeq\mathbf{Grp}.
\]

对含幺环 \(R\)，同样取有限秩自由左 \(R\)-模范畴的相反范畴 \(\mathcal L_R\)，规定
\[
\operatorname{Hom}_{\mathcal L_R}(n,m)
=\operatorname{Hom}_R(R^m,R^n).
\]
有限直和在模范畴中是余积，所以相反范畴中的 \(n+m\) 是积。态射 \(n\to1\) 由 \(1\) 的像 \(\sum_i r_ie_i\in R^n\) 决定，表示形式线性项 \(\sum_i r_ix_i\)。代入时，外层系数在左侧乘内层系数：
\[
\sum_j s_j\left(\sum_i r_{ji}x_i\right)
=\sum_i\left(\sum_j s_jr_{ji}\right)x_i.
\]
这一次序来自自由左模同态的复合，不能在非交换环上交换系数。

\ProseParagraphBreak
模型由零项、加法项、负元项及各 \(rx_1\) 定义底集合上的左 \(R\)-模结构。形式线性项的等式给出交换群公理、分配律与 \(r(sx)=(rs)x\)、\(1x=x\)。反向，一个左 \(R\)-模按有限线性组合求值，便给出模型；保持所有项的函数恰为左模同态。因此
\[
\operatorname{Mod}(\mathcal L_R)\simeq R\text{-}\mathbf{Mod}.
\]
群与模的例子说明了为何有限积必不可少：同一个变量可以被重复使用，也可以被遗漏，多元运算可以同时接受这些变量；投影及其积的泛性质把这些替换组织成态射。Lawvere 以这种方式研究代数理论的函子语义\cite[pp. 869--870]{Lawvere1963FunctorialSemantics}。


\begin{chaptersummary}
单子代数把形式运算变成求值，单位和结合律保证求值与归并相容。每个单子都有自由代数与遗忘伴随；原来的伴随能否由它恢复，则取决于比较函子是否为等价。Barr–Beck 判据把这个问题归结为反映同构以及正确处理遗忘后分裂的余等化子。

\ProseParagraphBreak

模的限制标量总具有单子性，不需要平坦性；忠实平坦下降则讨论扩张标量的余单子，平坦性保证等化子保留，忠实性保证原对象与态射没有被丢失。可表性和伴随存在性又有不同的困难：极限组织相容条件，解集把候选对象限制为集合大小，始对象的泛性质最终保证存在唯一性。

\ProseParagraphBreak
Kleisli 态射把一个对象送入另一个对象的形式表达式中，单子乘法定义替换的复合；它与自由单子代数之间的同态完全对应。Lawvere 理论进一步以有限积记录有限变量、投影、运算项与代入，保持积的集合模型给出运算的实际求值。群词与形式线性组合分别恢复群与左模，同时也恢复它们的同态。
\end{chaptersummary}

\BookExercises

\begin{exercise}\label{b4:ex:A-pointed-monad}
在集合范畴上令 \(TX=X\amalg\{*\}\)，单位为自然包含，乘法 \(T^2X\to TX\) 固定 \(X\) 中元素并把两个新元素都送到 \(*\)。验证单子公理，并识别其代数范畴。
\end{exercise}
\begin{solution}
在三次复合中，无论按哪一顺序合并，原来的 \(x\in X\) 都保持不变，任意一层的新元素都变成唯一的 \(*\)，故结合律成立。两条单位律也逐点成立。代数结构 \(a:X\amalg\{*\}\to X\) 由单位律迫使在 \(X\) 上为恒等，只剩指定 \(x_0=a(*)\)。第二条公理只要求所有新元素最后都送到 \(x_0\)，不增加条件。代数同态恰好是保持 \(x_0\) 的函数，所以代数范畴是带基点集合范畴；空集合不能带这样的代数结构。
\end{solution}

\begin{exercise}\label{b4:ex:A-nonempty-list}
把有限表单子中的空表删除，只保留非空有限表，仍用单元素表及连接定义单位和乘法。识别所得代数范畴，特别说明空底集合的情况，并指出它与幺半群情形的差别。
\end{exercise}
\begin{solution}
非空表的非空表仍能连接成非空表，两种连接顺序相同，单位也仍成立。若 \(a\) 为代数结构，定义 \(x\cdot y=a\bigl([x,y]\bigr)\)，对三元素表的两种分组应用代数公理，得到结合律。归纳给出所有非空表的求值都等于依次相乘，因此没有额外数据。反过来，结合乘法定义所有非空表的值，满足代数公理，同态恰好保持这个乘法。

这里 \(T\varnothing=\varnothing\)，所以空集合也有唯一的代数结构。因而完整的代数范畴允许空底集合上的结合乘法；限制到非空对象，才是通常要求底集合非空的半群范畴。构造中没有空表，所以没有被指定的单位元，也没有要求一个本来存在的单位元被同态保持。
\end{solution}

\begin{exercise}\label{b4:ex:A-cofree-coalgebra}
设 \(Q\) 为局部小范畴 \(\mathcal D\) 上的余单子。证明 \((QY,\delta_Y)\) 是余代数，并给出
\[
\operatorname{Hom}_{\mathcal D_Q}
\bigl((X,c),(QY,\delta_Y)\bigr)
\cong\operatorname{Hom}_{\mathcal D}(X,Y)
\]
的显式双射。这说明遗忘函子有什么伴随？
\end{exercise}
\begin{solution}
余单子的公理就是 \((QY,\delta_Y)\) 的余代数公理。把余代数同态 \(f:X\to QY\) 送到 \(\varepsilon_Yf\)；给定 \(h:X\to Y\)，逆映射为 \(Q(h)c\)。后者是余代数同态，因为
\[
\delta_YQ(h)c=Q^2(h)\delta_Xc
=Q^2(h)Q(c)c=Q\bigl(Q(h)c\bigr)c.
\]
余单位自然性及 \(\varepsilon_Xc=1_X\) 说明第一种复合为 \(h\)。反方向使用余代数同态条件：
\[
Q(\varepsilon_Yf)c
=Q(\varepsilon_Y)Q(f)c
=Q(\varepsilon_Y)\delta_Yf=f.
\]
公式与态射复合相容，故双射自然。遗忘函子因而有右伴随 \(Y\mapsto(QY,\delta_Y)\)，称为余自由余代数构造。
\end{solution}

\begin{exercise}\label{b4:ex:A-rationalization}
在交换群范畴上考虑单子 \(T(A)=\mathbb Q\otimes_{\mathbb Z}A\)。证明它幂等，并证明其代数范畴等价于 \(\mathbb Q\)-向量空间范畴。描述其固定对象在交换群语言中的性质。
\end{exercise}
\begin{solution}
单子乘法及其候选逆为
\[
\begin{aligned}
\mu_A\bigl(r\otimes(q\otimes a)\bigr)&=rq\otimes a,\\
\nu_A(q\otimes a)&=q\otimes(1\otimes a).
\end{aligned}
\]
定义 \(\nu_A\) 的映射对 \(q,a\) 双加性，且对每个整数 \(n\) 满足
\[
(nq)\otimes(1\otimes a)
=q\otimes\bigl(1\otimes(na)\bigr),
\]
故由张量积的泛性质诱导良定义同态。直接代入得 \(\mu_A\nu_A=1\)。

在 \(W=\mathbb Q\otimes_{\mathbb Z}(\mathbb Q\otimes_{\mathbb Z}A)\) 上，乘以任一非零整数 \(n\) 都可逆，其逆由第一个 \(\mathbb Q\) 因子乘以 \(1/n\) 诱导。把 \(q\) 写成 \(m/n\)，其中 \(m,n\) 为整数且 \(n\ne0\)，整数平衡关系给出
\[
\begin{aligned}
n\bigl(r\otimes(q\otimes a)\bigr)
&=r\otimes\bigl((nq)\otimes a\bigr)\\
&=mr\otimes(1\otimes a)
=n\bigl(rq\otimes(1\otimes a)\bigr).
\end{aligned}
\]
消去这个可逆同态，得到
\(r\otimes(q\otimes a)=rq\otimes(1\otimes a)\)，所以 \(\nu_A\mu_A=1\)。因此乘法为同构，单子幂等。

由\cref{b4:prop:A-restriction-scalars-monadic}，代数结构正是延伸原来整数作用的 \(\mathbb Q\)-作用。其底层交换群中，乘每个非零整数都是双射。反过来，若这些乘法均双射，则定义 \((a/b)x\) 为唯一满足 \(by=ax\) 的 \(y\)，通分与唯一性给出有理标量作用，故固定对象就是唯一可除交换群。幂等性说明这一作用一旦存在便唯一。
\end{solution}

\begin{exercise}\label{b4:ex:A-torsionfree-reflection}
对交换群 \(A\)，令 \(t(A)\) 为全部有限阶元素所成的子群。证明 \(A/t(A)\) 无挠，并证明 \(A\mapsto A/t(A)\) 左伴随于无挠交换群范畴的包含函子。说明相应单子为何幂等。
\end{exercise}
\begin{solution}
若 \(n\bigl(a+t(A)\bigr)=0\)，则 \(na\) 有限阶，所以存在 \(m>0\) 使 \(mna=0\)，即 \(a\in t(A)\)。因此商群无挠。到无挠群 \(B\) 的任何同态都把 \(t(A)\) 送到零，故唯一经过 \(A/t(A)\) 分解；这个商映射的泛性质对两变量自然，给出伴随。再次除去挠子群已经不改变无挠商，所以单子乘法为同构。此反射与有理化不同，例如 \(\mathbb Z\) 已经无挠，却不是唯一可除的。
\end{solution}

\begin{exercise}\label{b4:ex:A-em-limits}
设 \(T\) 为局部小范畴 \(\mathcal C\) 上的单子，\(D:J\to\mathcal C^T\) 为小图。若底层图 \(U^TD\) 在 \(\mathcal C\) 中有极限，证明该极限上有唯一代数结构使极限锥成为代数同态锥，且它给出 \(\mathcal C^T\) 中的极限。这里是否需要 \(T\) 保持极限？
\end{exercise}
\begin{solution}
记底层极限为 \(L\)，投影为 \(\pi_j:L\to X_j\)，各代数结构为 \(a_j:TX_j\to X_j\)。映射 \(a_jT(\pi_j):TL\to X_j\) 构成锥，因为图中的态射是代数同态。故存在唯一 \(a:TL\to L\)，使
\(\pi_ja=a_jT(\pi_j)\)。

在两条代数公理的两边分别复合所有 \(\pi_j\)，利用 \(a_j\) 的单位、结合律与自然性，得到相同结果。极限投影联合为单，故 \(a\) 满足两条公理。对任意代数锥，底层唯一的极限态射与代数结构相容，同样可在复合所有 \(\pi_j\) 后检验，所以它给出代数范畴中的极限。整个构造只向底层极限送入一个锥，未曾要求 \(TL\) 是某个极限，因此不需要 \(T\) 保持极限。
\end{solution}

\begin{exercise}\label{b4:ex:A-free-module-monad}
设 \(R\) 为含幺环，在集合范畴上令 \(TX\) 为以 \(X\) 为基的自由左 \(R\)-模的底集合，单位为 \(x\mapsto e_x\)，乘法为展开有限形式线性组合。证明这个单子的代数范畴等价于左 \(R\)-模范畴。
\end{exercise}
\begin{solution}
给定代数求值 \(a:TX\to X\)，定义
\[
0_X=a(0),\qquad
x+y=a(e_x+e_y),\qquad
rx=a(re_x).
\]
单位律给出 \(a(e_x)=x\)。代数结合律说明，对任意有限个 \(v_i\in TX\)，
\[
a\left(\sum_i r_i e_{a(v_i)}\right)
=a\left(\sum_i r_i v_i\right).
\]
因此在内层先作加法或标量乘法，与在自由模中先展开后求值相同。自由模中的加法交换、结合、零元与负元等式分别给出 \(X\) 中的相应等式；取内层 \(v=se_x\) 得 \(r(sx)=(rs)x\)，取 \(v=e_x+e_y\) 得分配律，单位元给出 \(1x=x\)。所以得到左 \(R\)-模。

反向，左 \(R\) 模给出求值映射 \(a\bigl(\sum_i r_i e_{x_i}\bigr)=\sum_i r_i x_i\)，它满足两条代数公理。上面的等式还说明，原来的 \(a\) 必取这种形式，所以构造互逆。与求值映射相容的函数恰好保持零、加法与标量乘法，也就是左模同态。
\end{solution}

\begin{exercise}\label{b4:ex:A-forgetful-not-full}
对带基点集合的遗忘函子，说明它忠实且反映同构，但不是全函子。给出一个底层为分裂满射而非同构的代数结构 \(a:TX\to X\)。
\end{exercise}
\begin{solution}
一个带基点映射由其底层函数确定，所以遗忘函子忠实；保持基点的双射，其逆也保持基点，所以反映同构。对同一个带基点二点集 \(\bigl(\{0,1\},0\bigr)\)，交换 \(0,1\) 的底层函数不保持基点，故函子不全。

取 \(X=\{0,1\}\)，令 \(a:X\amalg\{*\}\to X\) 固定 \(X\) 并把 \(*\) 送到 \(0\)。它与单位包含复合为恒等，因而是分裂满射；但 \(a(0)=a(*)\)，所以不是同构。单子代数的单位律只保证求值分裂满，不保证求值可逆。
\end{solution}

\begin{exercise}\label{b4:ex:A-monoid-presentation}
设 \(M\) 为幺半群，以每个 \(x\in M\) 为一个形式生成元 \([x]\) 构造自由幺半群。证明由关系
\[
[x][y]=[xy],\qquad []=[1_M]
\]
所得的商幺半群同构于 \(M\)。解释它与单子代数的自由对象展示的关系。
\end{exercise}
\begin{solution}
把一个词 \([x_1]\cdots[x_n]\) 送到 \(x_1\cdots x_n\)，空词送到 \(1_M\)，得到满同态，而且消去所列关系。反过来，关系允许把每个非空词依次缩成单个 \([x_1\cdots x_n]\)，把空词改成 \([1_M]\)。因此两个词只要求值相同，便在商中相同；若商中相同，求值当然相同。故核同余恰由这些关系生成，得到同构。

关系比较的正是“先在一层表内求值”与“先连接表再求值”。自由代数展示中的两支箭头分别执行这两种操作，余等化子把它们强制为同一个结果。
\end{solution}

\begin{exercise}\label{b4:ex:A-reflexive-not-split}
令 \(B=\{0,1,2\}\)、\(A=B\amalg\{a,b\}\)。定义 \(f,g:A\rightrightarrows B\)，在 \(B\) 上都为恒等，另令
\[
f(a)=0,\quad g(a)=1,\qquad f(b)=1,\quad g(b)=2.
\]
它们有共同截面 \(B\to A\)。求集合中的余等化子，并证明这个有序叉形图不能补成\cref{b4:def:A-split-coequalizer}中的分裂余等化子。
\end{exercise}
\begin{solution}
余等化关系包含 \(0\sim1\sim2\)，所以余等化子为单元素集合。若存在分裂，设截面 \(s\) 选择了 \(c\in B\)。条件 \(ft=1_B\)、\(gt=sq\) 要求：对每个 \(x\in B\)，都有一个 \(t(x)\) 满足 \(f\bigl(t(x)\bigr)=x\)、\(g\bigl(t(x)\bigr)=c\)。对 \(x=2\)，唯一的可能是 \(t(2)=2\)，所以 \(c=2\)。但没有元素同时满足 \(f=0,g=2\)，故 \(t(0)\) 不存在。共同截面使这对态射反身，却不足以给出所需的分裂余等化子。
\end{solution}

\begin{exercise}\label{b4:ex:A-abelian-forgetful-coequalizer}
证明交换群到集合的遗忘函子不保持全部余等化子。可以比较 \(\mathbb Z\) 上的恒等映射与取负映射。同时说明这为何不与该遗忘函子的单子性矛盾。
\end{exercise}
\begin{solution}
在交换群中，恒等与取负的差为乘二，所以余等化子为 \(\mathbb Z/2\mathbb Z\)。在集合中，生成的等价关系仅把 \(n\) 与 \(-n\) 识别，商集由 \(0,1,2,\ldots\) 标记，为无限集。因此底集合函子不保持这个余等化子。

由\cref{b4:ex:A-free-module-monad}取 \(R=\mathbb Z\)，该遗忘函子确有单子性。Barr–Beck 判据只要求保持底层可分裂的那些余等化子，不要求保持这一切平行态射的余等化子。此例也说明不能把定理中的限定删掉后再声称所得条件必要。
\end{solution}

\begin{exercise}\label{b4:ex:A-nonflat-monadic}
考虑 \(\mathbb Z\to\mathbb Z/2\mathbb Z\)。说明限制标量具有单子性，而扩张标量不保持单射。求对应单子的固定对象。
\end{exercise}
\begin{solution}
限制标量的单子性由\cref{b4:prop:A-restriction-scalars-monadic}得到，不需要平坦性。扩张标量把单射 \(\mathbb Z\xrightarrow{\times2}\mathbb Z\) 变成 \(\mathbb Z/2\mathbb Z\xrightarrow0\mathbb Z/2\mathbb Z\)，所以不保持单射。

单子为 \(T(A)=A/2A\)，再作用一次不变，故幂等。固定对象满足 \(2A=0\)，正是 \(\mathbb F_2\)-向量空间的底层交换群。其标量作用由加法唯一确定，这与幂等单子的代数结构唯一相符。
\end{solution}

\begin{exercise}\label{b4:ex:A-product-descent}
设 \(R\) 为交换环，\(S=R\times R\times R\)，底环映射为对角映射。写 \(S\)-模 \(M=M_1\oplus M_2\oplus M_3\)。仿照\cref{b4:prop:A-product-descent-model}，证明模下降数据等价于一族 \(R\)-模同构
\[
f_{ij}:M_j\longrightarrow M_i\qquad(1\le i,j\le3),
\]
满足 \(f_{ii}=1_{M_i}\) 及 \(f_{ij}f_{jk}=f_{ik}\)。具体求出恢复模，并举例说明仅要求 \(f_{ij}\)、\(f_{ji}\) 互逆，还不足以保证下降相容性。
\end{exercise}
\begin{solution}
令 \(e_i\in S\) 为第 \(i\) 个标准幂等元。在 \(S\otimes_RM\cong M^{\oplus3}\) 中，外层的第 \(i\) 个分量对应 \(e_i\)。\(S\)-线性迫使 \(\gamma\) 在 \(M_i\) 上的像落在这个分量，因而可唯一写成
\[
\gamma(m_1,m_2,m_3)
=\sum_{i=1}^3 e_i\otimes
\bigl(f_{1i}(m_i),f_{2i}(m_i),f_{3i}(m_i)\bigr).
\]
这里各 \(f_{ji}\) 为 \(R\)-线性映射。余单位从第 \(i\) 个外层分量取出其 \(M_i\) 部分，所以余单位律恰好给出 \(f_{ii}=1_{M_i}\)。在余结合律两边，取 \(S\otimes_RS\otimes_RM\) 的第 \((i,k)\) 个外层分量，再取其中的 \(M_j\) 分量，分别得到 \(f_{ji}(m_i)\) 和 \(f_{jk}f_{ki}(m_i)\)。因此余结合律恰好要求 \(f_{jk}f_{ki}=f_{ji}\)，即题中的相容等式。取 \(j=i\) 得 \(f_{ik}f_{ki}=1_{M_i}\)；交换 \(i,k\) 又得另一侧逆式，所以这些映射自动为同构。反过来，满足所列条件的映射族按上述公式定义下降数据。

等化条件 \(\gamma(m)=1\otimes m\) 要求 \(m_j=f_{ji}(m_i)\) 对所有 \(i,j\) 成立。由相容等式，只需指定 \(m_1\)，就唯一确定其余两分量。因此恢复模为
\[
V=\bigl\{\bigl(m,f_{21}(m),f_{31}(m)\bigr):m\in M_1\bigr\},
\]
投影到第一分量给出 \(V\cong M_1\)。

取 \(R=\mathbb Q\)、\(M_i=\mathbb Q\)，令 \(f_{21}=f_{12}=f_{32}=f_{23}=1\)，\(f_{31}\) 为乘二，\(f_{13}\) 为乘 \(1/2\)，且 \(f_{ii}=1\)。每对反向映射都互逆，却有 \(f_{32}f_{21}\ne f_{31}\)。所以三份模两两能够识别，还要检查这些识别在三份之间相容，才能得到下降数据。
\end{solution}

\begin{exercise}\label{b4:ex:A-rational-descent-loss}
对平坦但不忠实的映射 \(\mathbb Z\to\mathbb Q\)，证明每个 \(\mathbb Q\)-向量空间上的模下降数据唯一。比较原来的两个交换群 \(\mathbb Z\) 与 \(\mathbb Q\) 所产生的数据，说明恢复失去何种唯一性。
\end{exercise}
\begin{solution}
对 \(\mathbb Q\)-空间 \(M\)，余单位
\(\mathbb Q\otimes_{\mathbb Z}M\to M\)、\(q\otimes m\mapsto qm\) 为同构，其逆为 \(m\mapsto1\otimes m\)。由余单位条件，\(\gamma\) 必是这个逆，故唯一；它也满足余结合律。

\(\mathbb Q\otimes_{\mathbb Z}\mathbb Z\) 与
\(\mathbb Q\otimes_{\mathbb Z}\mathbb Q\) 都同构于 \(\mathbb Q\)，并带有同一个唯一数据，但 \(\mathbb Z\) 与 \(\mathbb Q\) 作为交换群不同构：前者非零元素不能无限次被二整除，后者可以。故扩张后的数据不能唯一决定原模。恢复等化子对这份数据给出的是 \(\mathbb Q\) 的底层交换群，并不能区分所有可能的原对象。
\end{solution}

\begin{exercise}\label{b4:ex:A-faithful-counit}
设 \(F\dashv U\) 为局部小范畴之间的伴随，余单位为 \(\varepsilon\)。证明 \(U\) 忠实，当且仅当每个 \(\varepsilon_Y:FUY\to Y\) 都为满态射。这与 \(U\) 反映同构有何区别？
\end{exercise}
\begin{solution}
若 \(U\) 忠实且 \(r\varepsilon_Y=s\varepsilon_Y\)，应用 \(U\) 后可消去具有右逆 \(\eta_{UY}\) 的 \(U\varepsilon_Y\)，得到 \(Ur=Us\)，再由忠实性得 \(r=s\)。所以 \(\varepsilon_Y\) 满。

反过来，若所有 \(\varepsilon_Y\) 满且 \(Ur=Us\)，余单位自然性给出
\[
r\varepsilon_Y=\varepsilon_ZF(Ur)
=\varepsilon_ZF(Us)=s\varepsilon_Y,
\]
消去后得到 \(r=s\)，故 \(U\) 忠实。底集合函子 \(\mathbf{Top}\to\mathbf{Set}\) 忠实，却不反映同构，因为连续双射不一定为同胚。因此这一判据不能替代 Barr–Beck 定理中的保守性。
\end{solution}

\begin{exercise}\label{b4:ex:A-nilpotent-not-representable}
在含幺交换环范畴上令 \(N(A)\) 为全部幂零元的集合。证明这个协变函子不可表。比较固定一个正整数 \(n\) 时的函子 \(N_n(A)=\{a:a^n=0\}\)。
\end{exercise}
\begin{solution}
若 \(N\) 由 \((B,u)\) 表示，则泛元素 \(u\in N(B)\) 满足某个 \(u^n=0\)。任何经环同态得到的像也都满足 \(n\) 次幂为零。但在 \(\mathbb Z[t]/(t^{n+1})\) 中，\(\bar t\) 幂零且 \(\bar t^n\neq0\)，不可能由 \(u\) 得到，与泛元素条件矛盾。所以 \(N\) 不可表。

固定 \(n\) 时，\(\mathbb Z[t]/(t^n)\) 连同 \(\bar t\) 表示 \(N_n\)。此处泛元素必须有一个确定的幂零指数，其所有像也就有同一个指数上界，因而不能产生指数任意大的幂零元。

\Cref{b4:ex:03-units-nilpotents}曾从保持积证明同一结论：在 \(A_r=k[t]/(t^r)\) 中，各 \(\bar t_r\) 都幂零，但元组 \((\bar t_r)_r\) 没有共同的幂零指数，所以不属于积环的幂零元集合。两种证明所发现的是同一个障碍：逐个存在幂零指数，不能保证有一个统一指数；前者把这个要求施加于泛元素，后者施加于积环的一个元素。
\end{solution}

\begin{exercise}\label{b4:ex:A-general-linear-representability}
设 \(R\) 为交换环、\(n\ge1\)。在交换 \(R\)-代数范畴上证明函子 \(A\mapsto\operatorname{GL}_n(A)\) 可表，并写出表示代数和泛元素。
\end{exercise}
\begin{solution}
令 \(X=(x_{ij})\) 为 \(n\times n\) 不定元矩阵，取
\[
B=R[x_{ij},t]/\left(t\det X-1\right).
\]
指定 \(R\)-代数同态 \(B\to A\)，等价于指定矩阵 \(M=(a_{ij})\) 及 \(b\in A\)，满足 \(b\det M=1\)。交换环上的伴随矩阵恒等式说明这等价于 \(M\) 可逆，而此时 \(b=(\det M)^{-1}\) 唯一。因此同态自然地对应于 \(\operatorname{GL}_n(A)\)，泛元素是 \((\bar x_{ij})\)。这里 \(A\) 不必是域，可逆性由行列式为单位来判断。
\end{solution}

\begin{exercise}\label{b4:ex:A-nonempty-sets-initial}
在所有非空集合及全部函数组成的范畴中，证明单元素集合为弱始对象，但不存在始对象。找出该范畴缺少的一个等化子，说明弱始对象集引理的完备性假设不能删去。
\end{exercise}
\begin{solution}
每个非空集合至少接收一个来自单元素集合的函数，所以单元素集合弱始。任一非空集合 \(X\) 到二元素集合至少有两个不同的常值函数，因此没有对象能为始对象。

在二元素集合上取恒等映射与交换两点的映射，它们在集合范畴中的等化子为空集。任何非空集合都不可能映入这个等化图，因为两映射没有相等的取值点。因此所述非空集合范畴没有这个等化子，不能应用弱始对象集引理。
\end{solution}

\begin{optionalexercise}\label{b4:ex:A-ordinal-solution-set}
考虑所有序数组成的范畴 \(\mathcal O\)，规定从 \(\alpha\) 到 \(\beta\) 存在唯一态射，当且仅当 \(\alpha\ge\beta\)。证明它局部小且有全部小极限，但没有始对象。由此给出一个保持小极限却没有左伴随的函子，并说明解集条件在哪里失败。
\end{optionalexercise}
\begin{solution}
每个同态集合为空或单元素，所以局部小。对任一小图，所有顶点序数形成一个集合，其上确界是该图的极限；空图的极限是 \(0\)。这是因为极限锥的顶点恰为所有这些序数的共同上界，最小共同上界给出唯一分解。

若有始对象 \(\alpha\)，就必须有 \(\alpha\ge\beta\) 对所有序数成立，与 \(\beta=\alpha+1\) 矛盾。到单对象单态射范畴 \(\mathbf1\) 的唯一函子保持小极限；如果它有左伴随，左伴随在唯一对象处的值就应为 \(\mathcal O\) 的始对象，故不存在。

对任意集合大小的候选序数族，取大于其上确界的序数 \(\beta\)，就没有任何候选对象到 \(\beta\) 的态射。因此不存在弱始对象集，解集条件恰好失败。
\end{solution}

\begin{exercise}\label{b4:ex:A-torsion-inclusion}
设 \(\mathbf{TorAb}\) 为挠交换群的满子范畴，\(J:\mathbf{TorAb}\to\mathbf{Ab}\) 为包含函子。证明 \(J\) 有右伴随，却没有左伴随。
\end{exercise}
\begin{solution}
到任意交换群 \(A\) 的一个挠群同态，其像都在 \(t(A)\) 中，所以
\[
\operatorname{Hom}_{\mathbf{Ab}}(JT,A)
\cong\operatorname{Hom}_{\mathbf{TorAb}}\bigl(T,t(A)\bigr).
\]
因此 \(t\) 是 \(J\) 的右伴随。

如果 \(J\) 有左伴随，则 \(J\) 本身是右伴随，应保持全部存在的积。取 \(T_n=\mathbb Z/2^n\mathbb Z\)。在 \(\mathbf{TorAb}\) 中，其积为
\(t(\prod_nT_n)\)：来自挠群的任意相容映射族在全积中的像仍为挠元素，故满足所需泛性质。但 \(\prod_nT_n\) 中的元素 \((\bar1)_n\) 有无限阶，所以这个挠子群不是全积。包含函子不保持该积，矛盾，故没有左伴随。把挠元取出来是余反射，而不能据此猜测一个反方向的“最大挠商”具有同样泛性质。
\end{solution}

\begin{exercise}\label{b4:ex:A-kleisli-partial-functions}
取集合范畴上的单子 \(TX=X\amalg\{*\}\)，单位与乘法如\cref{b4:ex:A-pointed-monad}。把其 Kleisli 范畴识别为集合与部分函数的范畴，写出复合，并证明其中的同构恰好为处处有定义的双射。
\end{exercise}
\begin{solution}
函数 \(f:X\to Y\amalg\{*\}\) 等价于指定子集 \(D_f\subseteq X\) 及函数 \(D_f\to Y\)：把取值为 \(*\) 的点理解为未定义。若 \(g:Y\to Z\amalg\{*\}\)，则 \(g\star f\) 在 \(x\) 处有定义，当且仅当 \(x\in D_f\) 且 \(f(x)\in D_g\)，这时其值为 \(g(f(x))\)。两种未定义情形均由单子乘法合并成 \(*\)。单位把每个点送到自身，正是处处有定义的恒等函数。

若 \(g\star f=1_X\)，则 \(f\) 在每个 \(x\) 处都必须有定义；若又 \(f\star g=1_Y\)，则 \(g\) 也处处有定义。这两个复合等式于是成为普通函数的互逆条件，故 \(f\) 是双射。反向，处处有定义的双射及其逆自然给出 Kleisli 同构。
\end{solution}

\begin{exercise}\label{b4:ex:A-kleisli-free-monoid-automorphisms}
在有限表单子的 Kleisli 范畴中，令 \(X=\{x_1,\ldots,x_n\}\)。证明 \(X\) 的自同构都来自字母的置换。等价地，证明有限生成自由幺半群的自同构必须把自由生成元置换起来。
\end{exercise}
\begin{solution}
Kleisli 自同构等价于自由幺半群 \(X^*\) 的自同构 \(\alpha\)，设逆为 \(\beta\)。当 \(n=0\) 时只有平凡幺半群，结论成立。若 \(n>0\)，则 \(\alpha(x_i)\) 不可能为空词，否则 \(\beta\alpha(x_i)\) 也为空，不等于 \(x_i\)。同理，每个 \(\beta(x_j)\) 都非空。

设 \(\alpha(x_i)\) 有 \(r\) 个字母。替换每个字母为非空的 \(\beta\) 像以后，词长至少为 \(r\)，而 \(\beta\alpha(x_i)=x_i\) 的词长为 \(1\)。故 \(r=1\)，每个生成元被送到一个生成元。单射性使这些像互不相同，有限性使它们穷尽全部生成元，所以给出一个置换。反过来，任意字母置换逐字延伸为自由幺半群自同构，其逆由逆置换给出。
\end{solution}

\begin{exercise}\label{b4:ex:A-kleisli-module-matrices}
设 \(R\) 为含幺环，\(T\) 为集合上的自由左 \(R\)-模单子。在其 Kleisli 范畴中，令 \(X\) 为单元素集合，并以 \(e\) 记 \(TX\) 的自由生成元。对 \(a\in R\)，令 \(f_a:X\to TX\) 将唯一元素送到 \(ae\)。计算 \(f_b\star f_a\)，识别 \(X\) 的自同态幺半群及自同构群，并与 \(\operatorname{End}_R({}_RR)\) 比较。最后取域 \(k\) 上的矩阵环 \(R=M_2(k)\)，用矩阵单位 \(a=E_{12}\)、\(b=E_{21}\) 说明复合次序不能忽略。
\end{exercise}
\begin{solution}
代入 \(f_b\) 后，外层系数 \(a\) 保留在左边，得到
\[
f_b\star f_a=f_{ab},\qquad 1_X^{\mathrm{Kl}}=f_1.
\]
因此按通常的“后复合乘在左边”约定，\(a\mapsto f_a\) 把自同态幺半群识别为 \(R\) 的乘法幺半群的相反幺半群。\(f_a\) 可逆，当且仅当存在 \(b\in R\) 满足 \(ab=ba=1\)，所以自同构群同构于 \((R^\times)^{\mathrm{op}}\)，也经 \(a\mapsto a^{-1}\) 同构于 \(R^\times\)。

\(f_a\) 延伸到自由左模 \(R\) 上的同态为右乘映射 \(\rho_a:r\mapsto ra\)，因为左线性迫使 \(\rho_a(r)=r\rho_a(1)\)。所有左模自同态都以此方式得到，且 \(\rho_b\circ\rho_a=\rho_{ab}\)，于是恢复 \(\operatorname{End}_R({}_RR)\cong R^{\mathrm{op}}\)。取所给矩阵单位，则
\[
f_b\star f_a=f_{E_{11}},\qquad
f_a\star f_b=f_{E_{22}}.
\]
这两映射不同，正好检验了正文矩阵计算中外层系数乘在内层系数左边的次序。
\end{solution}

\begin{exercise}\label{b4:ex:A-lawvere-commutative-monoids}
以 \(\mathbb N^n\) 为具有指定基的自由交换幺半群，构造它们组成的范畴的相反范畴 \(\mathcal L_{\mathrm{CMon}}\)。写出其态射的矩阵表示，证明它是 Lawvere 理论，且其集合模型范畴等价于交换幺半群范畴。
\end{exercise}
\begin{solution}
同态 \(\mathbb N^m\to\mathbb N^n\) 由 \(m\) 个基元素的像决定。把这些像的系数排列为 \(m\times n\) 矩阵 \(A\)，便得到理论中的态射 \(n\to m\)，其第 \(j\) 项为 \(\sum_i a_{ji}x_i\)。若随后态射 \(m\to r\) 的矩阵为 \(B\)，代入给出 \(BA\)。\(\mathbb N^{n+m}\) 是 \(\mathbb N^n,\mathbb N^m\) 的余积：到任一交换幺半群的同态只需分别指定两组基的像。因此相反范畴中 \(n+m\) 是积，而 \(0\) 为终对象，得到 Lawvere 理论。

模型中的零项 \(0\to1\) 与加法项 \((x_1+x_2):2\to1\) 给出一个零元和二元加法。形式和的结合、交换与零元等式使底集合成为交换幺半群。任何系数项的求值都被迫等于相应的重复加法。反过来，在交换幺半群中按重复加法求值每个非负整数系数项，代入等式由分配展开成立，故给出模型。保持这些项的函数正是保持零和加法的幺半群同态；积识别把两种构造互相恢复，于是得到所述等价。
\end{solution}

\begin{exercise}\label{b4:ex:A-lawvere-abelian-groups}
对群的 Lawvere 理论，在群词之外再要求 \(x_1x_2=x_2x_1\)，即把各自由群 \(F_n\) 换成其交换化 \(\mathbb Z^n\)。证明所得理论的态射为整数矩阵，其模型恰为交换群，并解释这条关系对模型态射的要求。
\end{exercise}
\begin{solution}
交换化后的每个 \(n\) 变量群词唯一写为 \(x_1^{a_1}\cdots x_n^{a_n}\)，其中 \(a_i\in\mathbb Z\)。同态 \(\mathbb Z^m\to\mathbb Z^n\) 由基元素的像决定，所以理论中的 \(n\to m\) 由 \(m\times n\) 整数矩阵给出；代入对应矩阵乘法。有限自由交换群的直和是余积，故相反范畴仍有指定有限积。

群理论的模型已经有群结构，新等式要求任意两个元素可交换。因此模型是交换群，各整数系数项按整数倍与加法求值。反向，任意交换群都满足这些等式，并由 \(\sum_i a_ix_i\) 给出保持积的模型。模型自然变换保持加法、零与负元，正是交换群同态；增加交换关系约束了对象的乘法，没有增加群同态之外的新态射数据。
\end{solution}
